\documentclass[11pt,a4paper]{article}

% ==================== TEMEL PAKETLER (pdflatex) ====================
\usepackage[T1]{fontenc}
\usepackage[utf8]{inputenc}
\usepackage{lmodern}
\usepackage[a4paper,margin=2cm,top=2.4cm,bottom=2.4cm]{geometry}
\usepackage{xcolor}
\usepackage{colortbl}
\usepackage{array}
\usepackage{booktabs}
\usepackage{longtable}
\usepackage{tabularx}
\usepackage{listings}
\usepackage{fancyhdr}
\usepackage{amsmath}
\usepackage{amssymb}
\usepackage{hyperref}

% ==================== RENK PALETİ ====================
\definecolor{anaMavi}{HTML}{132C4A}
\definecolor{vurguMavi}{HTML}{2E6DA4}
\definecolor{acikMavi}{HTML}{EAF2F9}
\definecolor{turuncu}{HTML}{D9701E}
\definecolor{yesil}{HTML}{2E7D32}
\definecolor{acikYesil}{HTML}{E8F5E9}
\definecolor{kirmizi}{HTML}{C0392B}
\definecolor{mor}{HTML}{6C3483}
\definecolor{acikMor}{HTML}{F3E9F7}
\definecolor{acikGri}{HTML}{F4F5F7}
\definecolor{ortaGri}{HTML}{6B7280}
\definecolor{koyuGri}{HTML}{2D2D2D}
\definecolor{kodArka}{HTML}{FBFBFA}

% ==================== HYPERREF ====================
\hypersetup{
  colorlinks=true,
  linkcolor={[HTML]{2E6DA4}},
  urlcolor={[HTML]{2E6DA4}},
  pdftitle={CMake ile Modern C/C++ Derleme Rehberi},
  pdfauthor={Derleme Sistemleri Rehberi},
  bookmarksnumbered=true
}

% ==================== BAŞLIK STİLLERİ (titlesec olmadan) ====================
\setcounter{secnumdepth}{0}
\newcounter{secc}
\newcounter{subsecc}[secc]
\newcounter{subsubsecc}[subsecc]
\newcommand{\gsec}[1]{\par\phantomsection\stepcounter{secc}%
  \setcounter{subsecc}{0}%
  \vspace{18pt}{\color{anaMavi}\Large\bfseries\thesecc.\hspace{0.5em}#1}\par
  \vspace{2pt}{\color{turuncu}\rule{\linewidth}{1.6pt}}\par\vspace{8pt}%
  \addcontentsline{toc}{section}{\protect\numberline{\thesecc}#1}\markboth{#1}{#1}}
\newcommand{\gsub}[1]{\par\phantomsection\stepcounter{subsecc}%
  \setcounter{subsubsecc}{0}%
  \vspace{11pt}{\color{vurguMavi}\large\bfseries\thesecc.\thesubsecc\hspace{0.5em}#1}\par\vspace{4pt}%
  \addcontentsline{toc}{subsection}{\protect\numberline{\thesecc.\thesubsecc}#1}}
\newcommand{\gsubsub}[1]{\par\phantomsection\stepcounter{subsubsecc}%
  \vspace{7pt}{\color{koyuGri}\normalsize\bfseries #1}\par\vspace{2pt}}

% Kısım (Part) başlığı
\newcommand{\gpart}[2]{\clearpage
  \begin{center}\vspace*{1.5cm}
  {\color{turuncu}\rule{0.7\linewidth}{1.5pt}}\\[10pt]
  {\color{ortaGri}\Large KISIM #1}\\[8pt]
  {\color{anaMavi}\Huge\bfseries #2}\\[10pt]
  {\color{turuncu}\rule{0.7\linewidth}{1.5pt}}
  \end{center}\vspace{1cm}}

% ==================== BAŞLIK / ALTBİLGİ ====================
\setlength{\headheight}{14pt}
\pagestyle{fancy}
\fancyhf{}
\renewcommand{\headrulewidth}{0.4pt}
\renewcommand{\footrulewidth}{0.4pt}
\fancyhead[L]{\small\color{ortaGri}CMake ile Modern C/C++}
\fancyhead[R]{\small\color{ortaGri}\leftmark}
\fancyfoot[C]{\small\color{ortaGri}\thepage}

% ==================== ORTAK LITERATE (Türkçe karakterler) ====================
% Her stilde tekrar kullanılacak dönüşümler.

% ==================== KOD LİSTELEME (CMake) ====================
\lstdefinestyle{cmakestyle}{
  basicstyle=\ttfamily\small\color{koyuGri},
  keywordstyle=\color{vurguMavi}\bfseries,
  keywordstyle=[2]\color{mor}\bfseries,
  keywordstyle=[3]\color{turuncu},
  commentstyle=\color{yesil}\itshape,
  stringstyle=\color{turuncu},
  numberstyle=\tiny\color{ortaGri},
  numbers=left,
  numbersep=8pt,
  backgroundcolor=\color{acikMavi!45},
  frame=single,
  framesep=6pt,
  rulecolor=\color{vurguMavi!40},
  breaklines=true,
  showstringspaces=false,
  tabsize=2,
  columns=fullflexible,
  keepspaces=true,
  extendedchars=true,
  morecomment=[l]{\#},
  sensitive=false,
  morekeywords={cmake_minimum_required,project,add_executable,add_library,
    add_subdirectory,target_link_libraries,target_include_directories,
    target_compile_features,target_compile_definitions,target_compile_options,
    target_sources,target_link_options,target_precompile_headers,
    set,unset,set_target_properties,get_target_property,set_property,
    get_property,option,list,string,file,math,separate_arguments,
    if,elseif,else,endif,foreach,endforeach,while,endwhile,
    function,endfunction,macro,endmacro,return,break,continue,
    include,include_directories,link_directories,link_libraries,
    add_definitions,add_compile_options,add_compile_definitions,
    find_package,find_library,find_path,find_program,find_file,
    install,export,configure_file,message,enable_testing,add_test,
    add_custom_command,add_custom_target,add_dependencies,source_group,
    FetchContent_Declare,FetchContent_MakeAvailable,FetchContent_GetProperties,
    FetchContent_Populate,ExternalProject_Add,write_basic_package_version_file,
    configure_package_config_file,install_targets,cmake_policy,
    cmake_parse_arguments,include_guard,define_property,mark_as_advanced},
  morekeywords=[2]{PUBLIC,PRIVATE,INTERFACE,STATIC,SHARED,MODULE,OBJECT,
    IMPORTED,ALIAS,REQUIRED,QUIET,EXACT,COMPONENTS,OPTIONAL_COMPONENTS,
    CACHE,FORCE,PARENT_SCOPE,GLOBAL,BOOL,STRING,PATH,FILEPATH,INTERNAL,
    AND,OR,NOT,STREQUAL,EQUAL,LESS,GREATER,MATCHES,VERSION_LESS,
    VERSION_GREATER,VERSION_EQUAL,DEFINED,EXISTS,IN,LISTS,ITEMS,RANGE},
  morekeywords=[3]{PROJECT_NAME,PROJECT_SOURCE_DIR,PROJECT_BINARY_DIR,
    CMAKE_SOURCE_DIR,CMAKE_BINARY_DIR,CMAKE_CURRENT_SOURCE_DIR,
    CMAKE_CURRENT_BINARY_DIR,CMAKE_CXX_STANDARD,CMAKE_C_STANDARD,
    CMAKE_BUILD_TYPE,CMAKE_INSTALL_PREFIX,CMAKE_RUNTIME_OUTPUT_DIRECTORY,
    CMAKE_CXX_STANDARD_REQUIRED,CMAKE_EXPORT_COMPILE_COMMANDS,
    CMAKE_MODULE_PATH,CMAKE_PREFIX_PATH,BUILD_SHARED_LIBS,
    CMAKE_CXX_COMPILER,CMAKE_C_COMPILER,CMAKE_SYSTEM_NAME},
  literate={ı}{{\i}}1{İ}{{\.I}}1{ş}{{\c s}}1{Ş}{{\c S}}1{ğ}{{\u g}}1{Ğ}{{\u G}}1{ç}{{\c c}}1{Ç}{{\c C}}1{ö}{{\"o}}1{Ö}{{\"O}}1{ü}{{\"u}}1{Ü}{{\"U}}1{—}{{-{}-}}2{–}{{-}}1{’}{{'}}1{“}{{"}}1{”}{{"}}1,
}

% ==================== KOD LİSTELEME (C++) ====================
\lstdefinestyle{cppstyle}{
  language=C++,
  basicstyle=\ttfamily\small\color{koyuGri},
  keywordstyle=\color{vurguMavi}\bfseries,
  commentstyle=\color{yesil}\itshape,
  stringstyle=\color{turuncu},
  numberstyle=\tiny\color{ortaGri},
  numbers=left,
  numbersep=8pt,
  backgroundcolor=\color{kodArka},
  frame=single,
  framesep=6pt,
  rulecolor=\color{ortaGri!40},
  breaklines=true,
  showstringspaces=false,
  tabsize=4,
  columns=fullflexible,
  keepspaces=true,
  extendedchars=true,
  literate={ı}{{\i}}1{İ}{{\.I}}1{ş}{{\c s}}1{Ş}{{\c S}}1{ğ}{{\u g}}1{Ğ}{{\u G}}1{ç}{{\c c}}1{Ç}{{\c C}}1{ö}{{\"o}}1{Ö}{{\"O}}1{ü}{{\"u}}1{Ü}{{\"U}}1{—}{{-{}-}}2{–}{{-}}1{’}{{'}}1{“}{{"}}1{”}{{"}}1,
  morekeywords={uint32_t,int32_t,uint8_t,size_t,nullptr,constexpr,noexcept,
    override,final,auto,std,string,vector,int64_t}
}

% ==================== KOD LİSTELEME (C) ====================
\lstdefinestyle{cstyle}{
  language=C,
  basicstyle=\ttfamily\small\color{koyuGri},
  keywordstyle=\color{kirmizi}\bfseries,
  commentstyle=\color{yesil}\itshape,
  stringstyle=\color{turuncu},
  numberstyle=\tiny\color{ortaGri},
  numbers=left,
  numbersep=8pt,
  backgroundcolor=\color{acikYesil!35},
  frame=single,
  framesep=6pt,
  rulecolor=\color{yesil!40},
  breaklines=true,
  showstringspaces=false,
  tabsize=4,
  columns=fullflexible,
  keepspaces=true,
  extendedchars=true,
  literate={ı}{{\i}}1{İ}{{\.I}}1{ş}{{\c s}}1{Ş}{{\c S}}1{ğ}{{\u g}}1{Ğ}{{\u G}}1{ç}{{\c c}}1{Ç}{{\c C}}1{ö}{{\"o}}1{Ö}{{\"O}}1{ü}{{\"u}}1{Ü}{{\"U}}1{—}{{-{}-}}2{–}{{-}}1{’}{{'}}1{“}{{"}}1{”}{{"}}1,
  morekeywords={uint32_t,int32_t,uint8_t,size_t,int64_t,bool,true,false}
}

% Kabuk/terminal stili
\lstdefinestyle{shstyle}{
  basicstyle=\ttfamily\small\color{koyuGri},
  backgroundcolor=\color{koyuGri!6},
  frame=single,framesep=6pt,rulecolor=\color{ortaGri!40},
  breaklines=true,showstringspaces=false,tabsize=4,columns=fullflexible,
  keepspaces=true,
  literate={ı}{{\i}}1{İ}{{\.I}}1{ş}{{\c s}}1{Ş}{{\c S}}1{ğ}{{\u g}}1{Ğ}{{\u G}}1{ç}{{\c c}}1{Ç}{{\c C}}1{ö}{{\"o}}1{Ö}{{\"O}}1{ü}{{\"u}}1{Ü}{{\"U}}1{—}{{-{}-}}2{–}{{-}}1,
}

\lstset{style=cmakestyle}

% ==================== ÖZEL KUTULAR (tcolorbox olmadan) ====================
\newcommand{\callout}[4]{\par\smallskip\noindent\begingroup
  \setlength{\fboxsep}{8pt}\setlength{\fboxrule}{0.8pt}%
  \fcolorbox{#1}{#2}{\parbox{\dimexpr\linewidth-2\fboxsep-2\fboxrule\relax}{%
    {\bfseries\color{#1}#3}\par\smallskip #4}}\endgroup\par\smallskip}
\newcommand{\bilgi}[1]{\callout{vurguMavi}{acikMavi}{Bilgi}{#1}}
\newcommand{\ipucu}[1]{\callout{yesil}{acikYesil}{İpucu}{#1}}
\newcommand{\uyari}[1]{\callout{kirmizi}{kirmizi!7}{Dikkat}{#1}}
\newcommand{\ornek}[2][Örnek]{\callout{ortaGri!90!black}{acikGri}{#1}{#2}}
\newcommand{\notk}[2][Not]{\callout{mor}{acikMor}{#1}{#2}}

% Mini proje banner'ı: her konunun sonundaki uygulamalı bölümü işaretler
\newcommand{\miniproj}[1]{\par\vspace{13pt}\noindent%
  {\setlength{\fboxsep}{7pt}\colorbox{turuncu}{\parbox{\dimexpr\linewidth-2\fboxsep\relax}{%
    \color{white}\bfseries MİNİ PROJE\hspace{0.6em}\textbullet\hspace{0.6em}#1}}}%
  \nopagebreak\par\vspace{7pt}}

% Yardımcı komutlar
\newcommand{\code}[1]{\texttt{\color{koyuGri}#1}}
\newcommand{\hd}[1]{\texttt{\color{koyuGri}\textless #1\textgreater}}
\renewcommand{\arraystretch}{1.35}

% ==================== BAŞLANGIÇ ====================
\begin{document}

% ---------- KAPAK ----------
\begin{titlepage}
\centering
\vspace*{1.4cm}
{\color{turuncu}\rule{\textwidth}{2pt}}\\[8pt]
{\color{anaMavi}\Huge\bfseries CMake ile Modern C/C++}\\[10pt]
{\color{vurguMavi}\LARGE Derleme Sistemleri, Bağımlılık Yönetimi ve Paketleme Rehberi}\\[8pt]
{\color{ortaGri}\large Tek Dosyadan Çok Modüllü Projeye, Sıfırdan İleri Seviyeye Kapsamlı Başvuru}\\[6pt]
{\color{turuncu}\rule{\textwidth}{2pt}}\\[1.4cm]

\begin{center}
\fcolorbox{vurguMavi}{acikMavi}{\parbox{0.86\textwidth}{\centering\large
Bu belge önce bir derleme sisteminin \textbf{neden gerektiğini} (derleyici,
bağlayıcı, \code{make}'in sınırları), sonra CMake'in \textbf{nasıl çalıştığını}
(configure/generate/build aşamaları, hedefler, özellik yayılımı) ve ardından
\textbf{gerçek projelerde} kütüphane yazmayı, bağımlılık çekmeyi, test ve
paketlemeyi \textbf{çok sayıda çalışan örnekle} anlatır. Örnekler hem
\textbf{C} hem \textbf{C++} için verilir ve modern (hedef-tabanlı) CMake
stiline uygundur.\vspace{4pt}}}
\end{center}

\vspace{0.7cm}
\begin{center}
\fcolorbox{ortaGri!50}{acikGri}{\parbox{0.86\textwidth}{\small
\textbf{Kapsam --- Kısım 1 (Temeller):} derleme nedir \textbullet\ önişlemci/derleyici/bağlayıcı
\textbullet\ \code{make}'in sorunları \textbullet\ CMake felsefesi \textbullet\ kurulum
\textbullet\ ilk C ve C++ projesi \textbullet\ out-of-source build \textbullet\ derleme akışı.\vspace{3pt}

\textbf{Kapsam --- Kısım 2 (Dil ve Hedefler):} değişkenler \& tipler \textbullet\ cache
\textbullet\ liste/string işlemleri \textbullet\ \code{add\_executable}/\code{add\_library}
\textbullet\ \code{target\_*} komutları \textbullet\ PUBLIC/PRIVATE/INTERFACE
\textbullet\ static/shared/interface/object kütüphaneler.\vspace{3pt}

\textbf{Kapsam --- Kısım 3 (Orta Seviye):} alt dizinler \textbullet\ proje düzeni
\textbullet\ \code{find\_package} \textbullet\ koşullar \& döngüler \textbullet\ fonksiyon/makro
\textbullet\ generator expression \textbullet\ derleme tipleri \textbullet\ sanitizer'lar.\vspace{3pt}

\textbf{Kapsam --- Kısım 4 (İleri):} FetchContent \textbullet\ ExternalProject
\textbullet\ install \& CPack \textbullet\ kendi config paketini yazmak
\textbullet\ CTest ile test \textbullet\ CMakePresets.json \textbullet\ toolchain \& çapraz derleme.\vspace{3pt}

\textbf{Kapsam --- Kısım 5 (Pratik):} tam çok-modüllü C++ projesi \textbullet\ tam C projesi
\textbullet\ en iyi uygulamalar \textbullet\ sık hatalar \& çözümleri \textbullet\ komut cheat sheet.\vspace{3pt}}}
\end{center}

\vfill
{\color{ortaGri}\large Hazırlanma Tarihi: 18 Temmuz 2026}\\[4pt]
{\color{ortaGri} CMake 3.20+ \textbullet\ C11/C17 \textbullet\ C++17/C++20 \textbullet\ GCC/Clang/MSVC \textbullet\ Ninja \& Make}
\vspace*{0.4cm}
\end{titlepage}

% ---------- İÇİNDEKİLER ----------
\tableofcontents
\thispagestyle{empty}
\clearpage

% #####################################################################
\gpart{1}{Temeller --- Derleme, make ve CMake}
% #####################################################################

% =====================================================================
\gsec{Bir Derleme Sistemi Neden Gerekir?}
% =====================================================================

CMake'i öğrenmenin ilk adımı, onun \emph{hangi soruna} çözüm getirdiğini
görmektir. Yazdığınız bir C ya da C++ programı olduğu gibi çalışmaz; kaynak kod
(\code{.c} ve \code{.cpp} dosyaları) önce bir dizi araçtan geçip makine koduna
dönüştürülür. Tek bir küçük dosyada bu adımları elle yürütmek kolaydır. Ne var
ki proje büyüdükçe --- onlarca dosya, farklı kütüphaneler, birden çok işletim
sistemi --- bu işi elle takip etmek olanaksız hale gelir. Derleme sistemleri de
tam olarak bu karmaşıklığı sizin yerinize yönetmek için vardır.

\gsub{Derleme Zinciri: Önişlemci, Derleyici, Bağlayıcı}

Bir kaynak dosyası çalıştırılabilir bir programa dönüşürken üç temel aşamadan
geçer. Bu aşamaları tanımak, ileride karşılaşacağınız hataları --- özellikle
``undefined reference'' türünden bağlayıcı hatalarını --- çözmenin anahtarıdır.

\begin{center}
\begin{tabularx}{\textwidth}{@{}>{\bfseries\color{anaMavi}}l >{\raggedright\arraybackslash}X >{\raggedright\arraybackslash}X@{}}
\toprule
\rowcolor{acikMavi}\color{anaMavi}\textbf{Aşama} & \color{anaMavi}\textbf{Ne yapar?} & \color{anaMavi}\textbf{Girdi $\to$ Çıktı}\\
\midrule
Önişlemci & \code{\#include}, \code{\#define}, \code{\#ifdef} gibi direktifleri işler; başlık dosyalarını metin olarak yerleştirir. & \code{.c/.cpp} $\to$ genişletilmiş kaynak\\
\rowcolor{acikGri} Derleyici & Genişletilmiş kaynağı makine koduna çevirir; her dosya için bir nesne (object) dosyası üretir. & \code{.cpp} $\to$ \code{.o}\\
Bağlayıcı (linker) & Tüm nesne dosyalarını ve kütüphaneleri birleştirir; sembol referanslarını çözer; tek bir çalıştırılabilir üretir. & \code{.o} + kütüphaneler $\to$ \code{a.out}\\
\bottomrule
\end{tabularx}
\end{center}

Örneğin \code{main.cpp} dosyasını elle derlemek şöyle görünür:

\begin{lstlisting}[style=shstyle]
# All in one step (preprocess + compile + link together)
g++ -std=c++17 main.cpp -o program

# Step by step: first the object file, then the linking stage
g++ -std=c++17 -c main.cpp -o main.o      # compile only: .cpp -> .o
g++ main.o -o program                      # link: .o -> executable

# Inspect the produced binary and its symbols
file program           # ELF 64-bit executable, dynamically linked, ...
nm main.o | head       # list symbols defined/referenced in the object file
\end{lstlisting}

\gsub{Elle Derlemenin Sınırları}

Tek dosya için yukarıdaki komut yeterli. Ama gerçek bir proje düşünün:

\begin{lstlisting}[style=shstyle]
# Compile every translation unit separately (-I adds a header search path)
g++ -std=c++17 -Iinclude -c src/main.cpp    -o build/main.o
g++ -std=c++17 -Iinclude -c src/matrix.cpp  -o build/matrix.o
g++ -std=c++17 -Iinclude -c src/vector.cpp  -o build/vector.o
g++ -std=c++17 -Iinclude -c src/parser.cpp  -o build/parser.o

# Link all object files together with the required libraries
# (-lpthread = POSIX threads, -lm = C math library)
g++ build/main.o build/matrix.o build/vector.o build/parser.o \
    -lpthread -lm -o program
\end{lstlisting}

Bu yaklaşımın sorunları:

\begin{itemize}
\item \textbf{Tekrar derleme:} Sadece \code{matrix.cpp} değişse bile, elle
her şeyi yeniden derlersiniz. Hangi dosyanın değiştiğini takip etmek gerekir.
\item \textbf{Taşınabilirlik yok:} Windows'ta \code{g++} yerine \code{cl.exe}
(MSVC), farklı bayraklar gerekir. Komutlarınız tek bir platforma bağlı kalır.
\item \textbf{Bağımlılık yönetimi:} Harici bir kütüphane (ör. Boost, SQLite)
kullanacaksanız onun yolunu, bayraklarını elle bulmak zorundasınız.
\item \textbf{Ölçeklenmez:} 200 dosyalık bir projede bu komutları yazmak ve
güncellemek imkansızdır.
\end{itemize}

\gsub{Make: İlk Çözüm ve Sınırları}

Bu soruna getirilen ilk büyük çözüm \code{make} oldu. \code{Makefile} adında
bir dosyaya ``hangi çıktı, hangi girdilerden ve nasıl üretilir'' kurallarını
yazarsınız; \code{make} da yalnızca \emph{değişen} dosyaları yeniden derler.

\begin{lstlisting}[style=shstyle]
# Variables reused throughout the recipes
CXX      = g++
CXXFLAGS = -std=c++17 -Iinclude
OBJ      = build/main.o build/matrix.o build/vector.o

# The final target depends on all object files; relink only when they change
program: $(OBJ)
	$(CXX) $(OBJ) -o program

# Pattern rule: how to build any build/X.o from src/X.cpp
# $< = the first prerequisite (the .cpp), $@ = the target (the .o)
build/%.o: src/%.cpp
	$(CXX) $(CXXFLAGS) -c $< -o $@
\end{lstlisting}

\code{make} güçlü bir araçtır, ama hâlâ platforma bağımlıdır: Windows'un yerel
derleyicisi \code{make} kullanmaz, Visual Studio kendi proje biçimini
(\code{.sln}/\code{.vcxproj}) bekler. Üstelik bağımlılıkları bulmayı, sürümleri
denetlemeyi, kurulum ve paketlemeyi de elle çözmeniz gerekir. CMake işte bu
boşluğu doldurmak için geliştirildi.

% =====================================================================
\gsec{CMake Nedir ve Nasıl Çalışır?}
% =====================================================================

CMake bir \emph{derleyici değildir}; kodu kendisi derlemez. O bir
\textbf{meta derleme sistemi}dir (build system generator): siz platformdan
bağımsız tek bir tarif (\code{CMakeLists.txt}) yazarsınız, CMake de bunu içinde
bulunduğunuz platformun \emph{yerel} derleme dosyalarına dönüştürür --- Linux'ta
Makefile ya da Ninja dosyalarına, Windows'ta Visual Studio projelerine, macOS'ta
Xcode projelerine. Böylece tek bir tarif her yerde çalışır.

\bilgi{Özet zihinsel model: \textbf{Sen} \code{CMakeLists.txt} yazarsın $\to$
\textbf{CMake} onu okuyup \code{build/} klasörüne Ninja/Make/VS dosyaları
\emph{üretir} (configure + generate) $\to$ \textbf{Ninja/Make} o dosyalarla
\code{g++}/\code{clang}/\code{cl}'i çağırıp kodu \emph{derler} (build).}

\gsub{İki Aşama: Configure/Generate ve Build}

CMake ile çalışmak her zaman iki mantıksal aşamadan oluşur:

\begin{center}
\begin{tabularx}{\textwidth}{@{}>{\bfseries\color{anaMavi}}l >{\raggedright\arraybackslash}X@{}}
\toprule
\rowcolor{acikMavi}\color{anaMavi}\textbf{Aşama} & \color{anaMavi}\textbf{Açıklama}\\
\midrule
Configure \& Generate & CMake \code{CMakeLists.txt} dosyalarını okur, derleyiciyi tespit eder, seçenekleri değerlendirir ve \code{build/} içine yerel derleme dosyalarını üretir. Komut: \code{cmake -S . -B build}\\
\rowcolor{acikGri} Build & Üretilen dosyalar kullanılarak asıl derleme yapılır (nesne dosyaları, kütüphaneler, çalıştırılabilirler). Komut: \code{cmake -{}-build build}\\
\bottomrule
\end{tabularx}
\end{center}

\gsub{Generator Kavramı}

CMake'in hangi yerel derleme sistemini üreteceğini ``generator'' (üretici)
belirler. En yaygın kullanılanlar şunlardır:

\begin{center}
\begin{tabularx}{\textwidth}{@{}>{\bfseries\color{anaMavi}}l >{\raggedright\arraybackslash}X@{}}
\toprule
\rowcolor{acikMavi}\color{anaMavi}\textbf{Generator} & \color{anaMavi}\textbf{Açıklama}\\
\midrule
Unix Makefiles & Klasik \code{make} dosyaları üretir (Linux/macOS varsayılanı).\\
\rowcolor{acikGri} Ninja & Çok hızlı, paralel derleme için ideal. Önerilir. \code{-G Ninja}\\
Visual Studio 17 2022 & Windows'ta \code{.sln} çözümü üretir.\\
\rowcolor{acikGri} Xcode & macOS'ta Xcode projesi üretir.\\
\bottomrule
\end{tabularx}
\end{center}

\begin{lstlisting}[style=shstyle]
cmake -S . -B build -G "Ninja"                    # use Ninja (recommended)
cmake -S . -B build -G "Unix Makefiles"           # use classic Make
cmake -S . -B build -G "Visual Studio 17 2022"    # generate a VS solution
cmake --help                                      # list all available generators
\end{lstlisting}

\ipucu{Ninja'yı öneririz: hem hızlıdır hem de tüm platformlarda aynı
davranır. Kurulum: Linux \code{sudo pacman -S ninja} / \code{apt install ninja-build},
macOS \code{brew install ninja}, Windows \code{winget install Ninja-build.Ninja}.}

% =====================================================================
\gsec{Kurulum ve İlk Proje}
% =====================================================================

\gsub{CMake Kurulumu}

\begin{lstlisting}[style=shstyle]
# Linux (Arch)
sudo pacman -S cmake ninja

# Linux (Debian/Ubuntu)
sudo apt install cmake ninja-build

# macOS
brew install cmake ninja

# Windows
winget install Kitware.CMake Ninja-build.Ninja

# Verify the installation (CMake 3.20 or newer is assumed in this guide)
cmake --version
ninja --version
\end{lstlisting}

\gsub{En Küçük C++ Projesi}

Bir klasör açalım ve iki dosya koyalım.

\begin{lstlisting}[style=cppstyle,title={\color{ortaGri}\ttfamily main.cpp}]
#include <iostream>

int main() {
    // std::cout is the standard output stream; std::endl flushes the buffer.
    std::cout << "Hello, CMake!" << std::endl;
    return 0;   // returning 0 from main() signals success to the OS
}
\end{lstlisting}

\begin{lstlisting}[title={\color{ortaGri}\ttfamily CMakeLists.txt}]
# Minimum required CMake version. Make this the first line of every project.
# It also decides which policy defaults ("modern behavior") are enabled.
cmake_minimum_required(VERSION 3.20)

# Project name, version and the languages used. CXX = C++, C = the C language.
# This call sets variables like PROJECT_NAME and PROJECT_VERSION for you.
project(HelloCMake VERSION 1.0 LANGUAGES CXX)

# Select the C++ standard used to compile all targets defined afterwards.
set(CMAKE_CXX_STANDARD 17)          # compile as C++17
set(CMAKE_CXX_STANDARD_REQUIRED ON) # fail instead of silently falling back
set(CMAKE_CXX_EXTENSIONS OFF)       # use -std=c++17, not the GNU -std=gnu++17

# Create an executable target: target name + its source files.
add_executable(hello main.cpp)
\end{lstlisting}

Şimdi derleyip çalıştıralım:

\begin{lstlisting}[style=shstyle]
cmake -S . -B build -G Ninja    # 1) configure + generate
cmake --build build            # 2) build (compiles main.cpp, links hello)
./build/hello                  # 3) run  ->  Hello, CMake!
\end{lstlisting}

\notk{\code{-S .} kaynak dizini (\code{CMakeLists.txt}'nin bulunduğu yer),
\code{-B build} ise derleme çıktılarının yazılacağı dizindir. Bu ayrım
\textbf{out-of-source build} denen iyi uygulamadır: kaynağınız temiz kalır,
\code{build/} klasörünü silerek her şeyi sıfırlayabilirsiniz.}

\gsub{En Küçük C Projesi}

Aynı yapı C için de neredeyse birebir aynıdır; tek fark \code{LANGUAGES C}
ve C standardıdır.

\begin{lstlisting}[style=cstyle,title={\color{ortaGri}\ttfamily main.c}]
#include <stdio.h>

int main(void) {
    /* printf writes formatted text to standard output; \n adds a newline. */
    printf("Hello, CMake (C)!\n");
    return 0;   /* 0 = success */
}
\end{lstlisting}

\begin{lstlisting}[title={\color{ortaGri}\ttfamily CMakeLists.txt}]
cmake_minimum_required(VERSION 3.20)
project(HelloC VERSION 1.0 LANGUAGES C)   # note: LANGUAGES C, not CXX

set(CMAKE_C_STANDARD 11)          # compile as C11
set(CMAKE_C_STANDARD_REQUIRED ON) # do not silently fall back to an older C
set(CMAKE_C_EXTENSIONS OFF)       # use -std=c11 instead of -std=gnu11

add_executable(hello_c main.c)
\end{lstlisting}

\begin{lstlisting}[style=shstyle]
cmake -S . -B build -G Ninja
cmake --build build
./build/hello_c    # -> Hello, CMake (C)!
\end{lstlisting}

\uyari{\code{build/} klasörünü \textbf{asla} sürüm kontrolüne (git) eklemeyin.
\code{.gitignore} dosyanıza \code{build/} satırını ekleyin. Bu klasör tümüyle
üretilir ve makineden makineye taşınamaz.}

\gsub{Faydalı Komut Satırı Bayrakları}

\begin{center}
\begin{tabularx}{\textwidth}{@{}>{\ttfamily\color{koyuGri}}l >{\raggedright\arraybackslash}X@{}}
\toprule
\rowcolor{acikMavi}\color{anaMavi}\textbf{Bayrak} & \color{anaMavi}\textbf{Ne yapar}\\
\midrule
-S <dir> & Kaynak dizini (CMakeLists.txt konumu).\\
\rowcolor{acikGri} -B <dir> & Derleme (binary) dizini.\\
-G <gen> & Generator seç (Ninja, "Unix Makefiles"...).\\
\rowcolor{acikGri} -D VAR=DEG & Bir cache değişkenine değer ata.\\
--build <dir> & Yapılandırılmış projeyi derle.\\
\rowcolor{acikGri} --build <dir> -j 8 & 8 iş parçacığıyla paralel derle.\\
--build <dir> -{}-target X & Sadece X hedefini derle.\\
\rowcolor{acikGri} --fresh & Cache'i yok sayıp baştan yapılandır (3.24+).\\
--install <dir> & Projeyi kur (install aşaması).\\
\bottomrule
\end{tabularx}
\end{center}

% #####################################################################
\gpart{2}{CMake Dili ve Hedefler}
% #####################################################################

% =====================================================================
\gsec{CMake Dilinin Temelleri}
% =====================================================================

\code{CMakeLists.txt} aslında küçük bir betik dilidir. Dosyadaki her satır bir
\emph{komut çağrısıdır}: \code{komut(argüman1 argüman2 ...)}. Bu dilin en
şaşırtıcı yanı şudur: \textbf{her şey birer metindir (string)}. Sayılar,
listeler, mantıksal değerler... hepsi arka planda metin olarak saklanır.

\gsub{Değişkenler ve set/unset}

\begin{lstlisting}
# Define a variable (its value is ALWAYS stored as text)
set(MY_NAME "Tolga")
set(MY_VERSION 42)

# Use a variable: ${...} "dereferences" it (substitutes its value)
message(STATUS "Name: ${MY_NAME}, Version: ${MY_VERSION}")

# A list is really a single string with ; as the element separator
set(SOURCES main.cpp util.cpp parser.cpp)   # stored as "main.cpp;util.cpp;parser.cpp"

# An undefined variable expands to an empty string (no error is raised)
message(STATUS "Undefined -> '${DOES_NOT_EXIST}'")   # -> Undefined -> ''

# Remove a variable
unset(MY_NAME)
\end{lstlisting}

\begin{lstlisting}[style=shstyle]
-- Name: Tolga, Version: 42
-- Undefined -> ''
\end{lstlisting}

\notk{\code{message(STATUS ...)} yapılandırma sırasında ekrana bilgi yazar ---
hata ayıklamanın en pratik yolu. Diğer seviyeler: \code{WARNING},
\code{FATAL\_ERROR} (yapılandırmayı durdurur), \code{DEBUG}.}

\gsub{Değişken Kapsamları (Scope)}

CMake'te bir değişken, varsayılan olarak tanımlandığı \emph{dizin} ve
\emph{fonksiyon} kapsamında yaşar. Bir alt dizine geçildiğinde
(\code{add\_subdirectory}) değişkenlerin bir kopyası o alt dizine aktarılır.
Bir değişikliği üst kapsama geri taşımak isterseniz \code{PARENT\_SCOPE}
anahtarını kullanırsınız.

\begin{lstlisting}
function(demo_scope)
    set(X "local")                  # visible only inside this function
    set(Y "outer" PARENT_SCOPE)     # written into the CALLER's scope
endfunction()

demo_scope()
message(STATUS "X='${X}' Y='${Y}'")   # -> X='' Y='outer'
\end{lstlisting}

\gsub{Cache Değişkenleri}

Sıradan değişkenler her yapılandırmada sıfırdan oluşturulur. \textbf{Cache
(önbellek) değişkenleri} ise \code{build/CMakeCache.txt} dosyasında kalıcı
olarak saklanır ve komut satırından \code{-D} ile değiştirilebilir. Kullanıcının
dışarıdan ayarlamasını istediğiniz seçenekler için bu tür değişkenleri kullanın.

\begin{lstlisting}
# A cache variable needs a type (STRING/BOOL/PATH/FILEPATH) and a description
set(MY_LIB_PATH "/usr/local" CACHE PATH "Library root directory")

# option() is a convenient shortcut for a boolean cache variable
option(BUILD_TESTS "Build the unit tests" ON)

# Using it
if(BUILD_TESTS)
    message(STATUS "Tests will be built")
endif()
\end{lstlisting}

\begin{lstlisting}[style=shstyle]
# Override cache values from the command line with -D:
cmake -S . -B build -D BUILD_TESTS=OFF -D MY_LIB_PATH=/opt/mylib

# Inspect or edit all cache variables interactively:
cmake -L build          # list cache variables and their values
ccmake build            # curses-based cache editor (if installed)
\end{lstlisting}

\bilgi{\code{option(X "..." ON)} aslında \code{set(X ON CACHE BOOL "...")}
kısaltmasıdır. Mantıksal değerler için \code{ON/OFF}, \code{TRUE/FALSE},
\code{1/0}, \code{YES/NO} kullanılabilir.}

\gsub{Liste İşlemleri}

\begin{lstlisting}
set(FILES a.cpp b.cpp c.cpp)

list(APPEND FILES d.cpp)             # append -> a;b;c;d
list(PREPEND FILES z.cpp)            # prepend -> z;a;b;c;d
list(REMOVE_ITEM FILES b.cpp)        # remove an element by value -> z;a;c;d
list(LENGTH FILES count)             # element count -> count = 4
list(GET FILES 0 first)              # index access -> first = z.cpp
list(FIND FILES c.cpp idx)           # index of an element -> idx = 2 (-1 if absent)
list(SORT FILES)                     # sort in place -> a;c;d;z
list(REVERSE FILES)                  # reverse in place -> z;d;c;a
list(REMOVE_DUPLICATES FILES)        # drop repeated elements, keep first
list(JOIN FILES ", " joined)         # list -> string: "z.cpp, d.cpp, c.cpp, a.cpp"
\end{lstlisting}

\gsub{String İşlemleri}

\begin{lstlisting}
string(TOUPPER "hello" upper)                # upper  = "HELLO"
string(TOLOWER "CMake" lower)                # lower  = "cmake"
string(LENGTH "cmake" len)                   # len    = 5
string(REPLACE "a" "@" out "cmake")          # out    = "cm@ke"
string(SUBSTRING "cmake" 0 2 head)           # head   = "cm"  (start, length)
string(FIND "hello.cpp" ".cpp" pos)          # pos    = 5  (index, -1 if absent)
string(REGEX MATCH "[0-9]+" num "abc123")    # num    = "123"  (first match)
string(REGEX REPLACE "[0-9]" "#" masked "a1b2")  # masked = "a#b#"
\end{lstlisting}

\gsub{Koşullar}

\begin{lstlisting}
if(WIN32)
    message(STATUS "We are on Windows")
elseif(APPLE)
    message(STATUS "We are on macOS")
elseif(UNIX)
    message(STATUS "We are on Linux/Unix")
endif()

# Comparison operators
if(CMAKE_CXX_COMPILER_VERSION VERSION_GREATER_EQUAL "11.0")  # version compare
endif()
if(name STREQUAL "tolga")                  # string equality
endif()
if(x EQUAL 42)                             # numeric equality
endif()
if(DEFINED MY_VAR)                         # is the variable defined?
endif()
if(EXISTS "${CMAKE_SOURCE_DIR}/config.h")  # does the file/directory exist?
endif()
if(NOT BUILD_TESTS AND USE_GUI)            # boolean operators: NOT / AND / OR
endif()
if("release" IN_LIST ALLOWED_TYPES)        # membership test on a list variable
endif()
\end{lstlisting}

\gsub{Döngüler}

\begin{lstlisting}
# Iterate over the elements of a list
set(MODULES core net ui)
foreach(mod IN LISTS MODULES)
    message(STATUS "Module: ${mod}")
endforeach()

# Iterate over several items directly (IN ITEMS)
foreach(name IN ITEMS alice bob carol)
    message(STATUS "User: ${name}")
endforeach()

# Numeric range (inclusive): 1,2,3,4,5
foreach(i RANGE 1 5)
    message(STATUS "i = ${i}")
endforeach()

# while loop with arithmetic via math(EXPR ...)
set(n 3)
while(n GREATER 0)
    message(STATUS "n = ${n}")
    math(EXPR n "${n} - 1")   # evaluate an integer expression
endwhile()
\end{lstlisting}

% =====================================================================
\gsec{Hedefler (Targets): Modern CMake'in Kalbi}
% =====================================================================

Modern CMake'in en temel kavramı \textbf{hedeftir} (target). Hedef, derlenecek
bir şeydir: bir çalıştırılabilir program ya da bir kütüphane. Eski CMake,
\code{include\_directories()} veya \code{add\_definitions()} gibi \emph{global}
komutlara dayanırdı; bunlar tüm projeyi birden etkilediği için büyük projelerde
işler kolayca içinden çıkılmaz hale gelirdi. Modern yaklaşım ise her ayarı
\textbf{ilgili hedefe bağlar}: ``bu hedef şu başlıklara, şu bayraklara ve şu
kütüphanelere ihtiyaç duyar'' dersiniz.

\bilgi{Altın kural: \textbf{Global komutlardan kaçın}
(\code{include\_directories}, \code{link\_libraries}, \code{add\_definitions}).
Bunun yerine \code{target\_*} komutlarını kullan. Her ayar, ait olduğu hedefe
bağlansın.}

\gsub{add\_executable ve add\_library}

\begin{lstlisting}
# Executable target
add_executable(app main.cpp app.cpp)

# Static library (.a / .lib) -- compiled into and embedded in the consumer
add_library(math STATIC matrix.cpp vector.cpp)

# Shared library (.so / .dll / .dylib) -- loaded at run time
add_library(math_shared SHARED matrix.cpp vector.cpp)

# If you omit the type, the BUILD_SHARED_LIBS variable decides
# (defaults to STATIC when BUILD_SHARED_LIBS is unset)
add_library(math_auto matrix.cpp vector.cpp)
\end{lstlisting}

\gsub{PUBLIC, PRIVATE, INTERFACE --- Yayılım (Propagation)}

\code{target\_*} komutlarındaki bu üç anahtar sözcük, bir ayarın \emph{kimi}
etkileyeceğini belirler. Bu üçünü kavramak, modern CMake'i kavramanın ta
kendisidir.

\begin{center}
\begin{tabularx}{\textwidth}{@{}>{\bfseries\color{anaMavi}}l >{\raggedright\arraybackslash}X@{}}
\toprule
\rowcolor{acikMavi}\color{anaMavi}\textbf{Anahtar} & \color{anaMavi}\textbf{Anlamı}\\
\midrule
PRIVATE & Sadece \emph{bu hedefin kendisi} kullanır. Bu hedefe bağlanan başkaları etkilenmez. (Uygulama detayı.)\\
\rowcolor{acikGri} INTERFACE & Bu hedefin kendisi kullanmaz ama \emph{ona bağlananlar} kullanır. (Sadece başlıklı --- header-only --- kütüphaneler için tipik.)\\
PUBLIC & Hem bu hedef hem de \emph{ona bağlananlar} kullanır. (PRIVATE + INTERFACE.)\\
\bottomrule
\end{tabularx}
\end{center}

\ornek[Sezgi]{Bir kütüphane, başlığında (\code{.h}) bir tipi
\code{\#include} ediyorsa, o başlığı kullanan \emph{herkes} de o include
yoluna ihtiyaç duyar $\to$ \textbf{PUBLIC}. Sadece \code{.cpp} içinde,
gizli bir şekilde kullanıyorsa $\to$ \textbf{PRIVATE}.}

\gsub{target\_include\_directories}

Bir hedefin başlık dosyalarını nerede arayacağını söyler.

\begin{lstlisting}
add_library(math STATIC src/matrix.cpp src/vector.cpp)

target_include_directories(math
    PUBLIC
        ${CMAKE_CURRENT_SOURCE_DIR}/include   # consumers of math.h see this too
    PRIVATE
        ${CMAKE_CURRENT_SOURCE_DIR}/src       # only math's own source files
)
\end{lstlisting}

\gsub{target\_link\_libraries}

Bir hedefi başka bir hedefe ya da kütüphaneye bağlar. Bu komutun asıl gücü
şuradadır: bağlandığınız hedefin \code{PUBLIC} ve \code{INTERFACE} özellikleri
--- include yolları, tanımlar, bayraklar --- \emph{kendiliğinden} sizin
hedefinize de geçer.

\begin{lstlisting}
add_library(math STATIC src/matrix.cpp)
target_include_directories(math PUBLIC include)

add_executable(app main.cpp)
# app AUTOMATICALLY inherits math's PUBLIC include path!
target_link_libraries(app PRIVATE math)
\end{lstlisting}

\ipucu{\code{target\_link\_libraries(app PRIVATE math)} yazdığınızda
\code{math/include} yolunu ayrıca \code{app}'a eklemeniz gerekmez.
Bu ``kullanım gereksinimlerinin'' (usage requirements) otomatik yayılımı,
modern CMake'in en güçlü özelliğidir.}

\gsub{target\_compile\_features --- Standart Belirlemenin Doğru Yolu}

\code{set(CMAKE\_CXX\_STANDARD 17)} tüm projeyi birden etkiler. Standardı hedef
hedef belirlemek ve bu gereksinimi bağlanan hedeflere de yaymak için
\code{target\_compile\_features} tercih edilir.

\begin{lstlisting}
add_library(math STATIC src/matrix.cpp)

# This library requires C++17; anything that links it also gets at least C++17.
target_compile_features(math PUBLIC cxx_std_17)

# For C:
# target_compile_features(mylib PUBLIC c_std_11)
\end{lstlisting}

\gsub{target\_compile\_definitions ve target\_compile\_options}

\begin{lstlisting}
add_executable(app main.cpp)

# Preprocessor macros (equivalent to -DNAME=VALUE on the command line)
target_compile_definitions(app PRIVATE
    APP_VERSION="1.2.0"      # in C++: #define APP_VERSION "1.2.0"
    ENABLE_LOGGING           # value-less macro, tested with #ifdef ENABLE_LOGGING
)

# Compiler flags (use with care -- these are not portable!)
target_compile_options(app PRIVATE
    -Wall -Wextra -Wpedantic
)
\end{lstlisting}

\uyari{Derleyici bayrakları taşınabilir değildir: \code{-Wall} GCC/Clang'te
çalışır ama MSVC'de \code{/W4} gerekir. Bu tür bayrakları generator expression
(Kısım 3) ile derleyiciye göre koşullamak gerekir.}

\gsub{Bir Hedefin Tüm Özelliklerini Bir Arada}

\begin{lstlisting}
add_library(graphics STATIC
    src/renderer.cpp
    src/texture.cpp
    src/shader.cpp
)

target_include_directories(graphics
    PUBLIC  include            # for consumers that include graphics.h
    PRIVATE src                # internal implementation details
)

target_compile_features(graphics PUBLIC cxx_std_17)

target_compile_definitions(graphics
    PUBLIC  GRAPHICS_VERSION="2.0"
    PRIVATE GRAPHICS_INTERNAL_BUILD
)

target_link_libraries(graphics
    PUBLIC  math               # if math types appear in graphics' public headers
    PRIVATE pthread            # used only inside the implementation
)
\end{lstlisting}

\miniproj{Yayılımı canlı görmek: uygulama hiçbir yol/makro yazmadan derleniyor}

Aşağıdaki projede \code{app}, ne başlık yolu ne de bir makro tanımlar; hepsini
\code{greet} kütüphanesinin \code{PUBLIC} özelliklerinden \emph{otomatik}
devralır. İşte modern CMake'in tüm mantığı budur.

\begin{lstlisting}[style=shstyle]
targets-demo/
|-- CMakeLists.txt
|-- include/greet.h
|-- src/greet.cpp
`-- app/main.cpp
\end{lstlisting}

\begin{lstlisting}[style=cppstyle,title={\color{ortaGri}\ttfamily include/greet.h}]
#pragma once
#include <string>
namespace greet { std::string hello(); }
\end{lstlisting}

\begin{lstlisting}[style=cppstyle,title={\color{ortaGri}\ttfamily src/greet.cpp}]
#include "greet.h"
namespace greet { std::string hello() { return "Merhaba, hedefler!"; } }
\end{lstlisting}

\begin{lstlisting}[style=cppstyle,title={\color{ortaGri}\ttfamily app/main.cpp}]
#include <iostream>
#include "greet.h"                 // include yolu greet'ten otomatik geldi
int main() {
    std::cout << greet::hello() << "\n";
    std::cout << "Surum: " << GREET_VERSION << "\n";  // makro da greet'ten geldi
}
\end{lstlisting}

\begin{lstlisting}[title={\color{ortaGri}\ttfamily CMakeLists.txt}]
cmake_minimum_required(VERSION 3.20)
project(TargetsDemo LANGUAGES CXX)

add_library(greet STATIC src/greet.cpp)
target_include_directories(greet PUBLIC include)               # app da gorecek
target_compile_definitions(greet PUBLIC GREET_VERSION="1.0")   # app da gorecek
target_compile_features(greet PUBLIC cxx_std_17)               # app da gorecek

add_executable(app app/main.cpp)
target_link_libraries(app PRIVATE greet)   # ustteki 3 PUBLIC ozellik app'e gecer
# DIKKAT: app icin ne include yolu ne de makro yazdik -- gerek kalmadi!
\end{lstlisting}

\begin{lstlisting}[style=shstyle]
cmake -S . -B build -G Ninja && cmake --build build
./build/app
#   Merhaba, hedefler!
#   Surum: 1.0
\end{lstlisting}

% =====================================================================
\gsec{Kütüphane Türleri}
% =====================================================================

\begin{center}
\begin{tabularx}{\textwidth}{@{}>{\bfseries\color{anaMavi}}l >{\raggedright\arraybackslash}X@{}}
\toprule
\rowcolor{acikMavi}\color{anaMavi}\textbf{Tür} & \color{anaMavi}\textbf{Açıklama}\\
\midrule
STATIC & \code{.a}/\code{.lib} --- kod, kullanan programa gömülür. Bağımlılık dağıtmak kolaydır.\\
\rowcolor{acikGri} SHARED & \code{.so}/\code{.dll}/\code{.dylib} --- çalışma anında yüklenir. Bellek paylaşımı, ayrı güncelleme.\\
MODULE & Çalışma anında \code{dlopen} ile açılan eklenti (plugin). Bağlanmaz.\\
\rowcolor{acikGri} INTERFACE & Derlenecek kaynağı yok; sadece kullanım gereksinimleri (header-only kütüphaneler).\\
OBJECT & Nesne dosyaları (\code{.o}) üretir, arşivlenmez; başka hedeflere gömülmek için.\\
\bottomrule
\end{tabularx}
\end{center}

Bu türlerden en sık kullandığınız ikisi \textbf{statik} ve \textbf{paylaşımlı
(dinamik)} kütüphanelerdir. İkisinin de kaynak kodu ve \code{add\_library}
çağrısı neredeyse birebir aynıdır; aradaki fark, kütüphanenin makine kodunun
\emph{ne zaman} ve \emph{nasıl} programa dahil olduğudur. Aşağıda ikisini de
sıfırdan, çalışan birer projeyle kuruyoruz. Her iki proje de \emph{aynı} kaynak
dosyalarını kullanır; yalnızca \code{CMakeLists.txt} içindeki tek kelime değişir.

\gsub{Statik Kütüphane (STATIC): Adım Adım}

Statik kütüphane derlendiğinde bir \code{.a} dosyası (Windows'ta \code{.lib})
üretir. Bu dosyadaki makine kodu, kütüphaneyi kullanan programa \emph{bağlama
anında} kopyalanır. Sonuçta ortaya tek ve kendi kendine yeten bir
çalıştırılabilir çıkar: dağıtırken yanında ayrıca kütüphane taşımanız gerekmez.
Bedeli ise, ikili dosyanın biraz büyümesi ve kütüphaneyi güncellediğinizde
programı yeniden bağlamanız gerekmesidir.

\miniproj{Statik \code{geo} kütüphanesi + onu kullanan uygulama}

Şu klasör yapısını kuralım:

\begin{lstlisting}[style=shstyle]
geo-static/
|-- CMakeLists.txt
|-- include/geo.h          # kutuphanenin genel arayuzu (baslik)
|-- src/geo.cpp            # kutuphanenin govdesi (kaynak)
`-- app/main.cpp           # kutuphaneyi kullanan uygulama
\end{lstlisting}

\begin{lstlisting}[style=cppstyle,title={\color{ortaGri}\ttfamily include/geo.h}]
#pragma once
namespace geo {
    double circle_area(double r);        // daire alani
    double rect_area(double w, double h); // dikdortgen alani
}
\end{lstlisting}

\begin{lstlisting}[style=cppstyle,title={\color{ortaGri}\ttfamily src/geo.cpp}]
#include "geo.h"
namespace geo {
    static const double PI = 3.14159265358979;
    double circle_area(double r)         { return PI * r * r; }
    double rect_area(double w, double h) { return w * h; }
}
\end{lstlisting}

\begin{lstlisting}[style=cppstyle,title={\color{ortaGri}\ttfamily app/main.cpp}]
#include <iostream>
#include "geo.h"

int main() {
    std::cout << "Daire (r=2)    : " << geo::circle_area(2) << "\n";
    std::cout << "Dikdortgen 3x4 : " << geo::rect_area(3, 4) << "\n";
}
\end{lstlisting}

\begin{lstlisting}[title={\color{ortaGri}\ttfamily CMakeLists.txt}]
cmake_minimum_required(VERSION 3.20)
project(GeoStatic VERSION 1.0 LANGUAGES CXX)

set(CMAKE_CXX_STANDARD 17)
set(CMAKE_CXX_STANDARD_REQUIRED ON)

# 1) STATIC kutuphane: derlenince libgeo.a (arsiv) uretir
add_library(geo STATIC src/geo.cpp)

# 2) geo.h'yi kullanan herkes include/ dizinini gorebilsin (PUBLIC)
target_include_directories(geo PUBLIC ${CMAKE_CURRENT_SOURCE_DIR}/include)

# 3) Uygulamayi tanimla ve kutuphaneye bagla
add_executable(app app/main.cpp)
target_link_libraries(app PRIVATE geo)   # geo'nun kodu app'in icine gomulur
\end{lstlisting}

\begin{lstlisting}[style=shstyle]
cmake -S . -B build -G Ninja
cmake --build build
./build/app
#   Daire (r=2)    : 12.5664
#   Dikdortgen 3x4 : 12

ls build/libgeo.a        # <- uretilen statik kutuphane (arsiv dosyasi)
ldd build/app | grep geo # <- CIKTI YOK: geo'nun kodu app'e gomuldu,
                         #    calisma aninda ayri bir dosyaya ihtiyac yok
\end{lstlisting}

\bilgi{\code{ldd} komutunun çıktısında \code{geo}'yu \emph{görmezsiniz}; çünkü
statik kütüphanede kod, çalıştırılabilirin \emph{içine} kopyalanmıştır.
\code{./build/app} dosyasını başka bir makineye kopyalasanız, \code{libgeo.a}
olmadan da çalışır.}

\gsub{Dinamik (Paylaşımlı) Kütüphane (SHARED): Adım Adım}

Paylaşımlı kütüphane; Linux'ta \code{.so}, Windows'ta \code{.dll}, macOS'ta
\code{.dylib} dosyası üretir. Bu kez kod programa kopyalanmaz, ayrı bir dosyada
kalır ve \emph{çalışma anında} yüklenir. Böylece aynı kütüphaneyi birçok program
paylaşabilir ve kütüphaneyi, programları yeniden derlemeden güncelleyebilirsiniz.
Karşılığında bir sorumluluk gelir: program çalışırken bu \code{.so} dosyasını
\emph{bulabilmelidir}; aksi halde ``cannot open shared object file'' hatası
alırsınız.

\miniproj{Aynı \code{geo} projesi, bu kez paylaşımlı kütüphaneyle}

Kaynak dosyaların (\code{geo.h}, \code{geo.cpp}, \code{main.cpp}) hepsi
\textbf{birebir aynı} kalır. Değişen tek şey \code{CMakeLists.txt}:

\begin{lstlisting}[title={\color{ortaGri}\ttfamily CMakeLists.txt}]
cmake_minimum_required(VERSION 3.20)
project(GeoShared VERSION 1.0 LANGUAGES CXX)

set(CMAKE_CXX_STANDARD 17)
set(CMAKE_CXX_STANDARD_REQUIRED ON)

# TEK FARK: STATIC yerine SHARED -> libgeo.so uretir
add_library(geo SHARED src/geo.cpp)
target_include_directories(geo PUBLIC ${CMAKE_CURRENT_SOURCE_DIR}/include)

# Iyi aliskanlik: kutuphaneye surum ver -> libgeo.so -> libgeo.so.1 -> .so.1.0
set_target_properties(geo PROPERTIES VERSION 1.0 SOVERSION 1)

add_executable(app app/main.cpp)
target_link_libraries(app PRIVATE geo)   # geo bu sefer app'e GOMULMEZ, disarida kalir
\end{lstlisting}

\begin{lstlisting}[style=shstyle]
cmake -S . -B build -G Ninja
cmake --build build
./build/app
#   Daire (r=2)    : 12.5664
#   Dikdortgen 3x4 : 12

ls build/libgeo.so*        # libgeo.so  libgeo.so.1  libgeo.so.1.0
ldd build/app | grep geo   # libgeo.so.1 => .../build/libgeo.so.1
                           # <- calisma aninda GEREKLI: kod ayri dosyada
\end{lstlisting}

\ipucu{Yukarıda \code{./build/app} sorunsuz çalıştı; çünkü CMake, derleme
ağacındaki kütüphanenin yolunu çalıştırılabilirin içine (RPATH olarak) otomatik
gömer. Kütüphaneyi \code{install} ile başka bir yere kurup taşırsanız,
çalışma anında bulunması için ya kurulum RPATH'i ayarlanmalı ya da
\code{LD\_LIBRARY\_PATH} kütüphanenin dizinini içermelidir.}

\notk[Hangisini seçmeli?]{\textbf{Statik}: dağıtımı kolaydır (tek dosya),
çalışma anında bağımlılık yoktur; ama ikili büyür ve güncelleme için yeniden
bağlama gerekir. \textbf{Paylaşımlı}: bellek ve disk paylaşımı sağlar, ayrı
güncellenebilir; ama \code{.so} çalışma anında bulunmalıdır. Kararsızsanız
statikle başlayın. \code{add\_library}'de tür yazmazsanız seçimi
\code{BUILD\_SHARED\_LIBS} değişkeni yapar: \code{-D BUILD\_SHARED\_LIBS=ON} ile
projedeki tüm kütüphaneleri tek hamlede paylaşımlıya çevirebilirsiniz.}

\gsub{Header-Only Kütüphane (INTERFACE)}

Yalnızca başlık dosyalarından oluşan bir kütüphaneyi (birçok modern C++
kütüphanesi böyledir) \code{INTERFACE} olarak tanımlarsınız; ortada derlenecek
bir \code{.cpp} yoktur.

\begin{lstlisting}
# Only header files (.hpp) under include/ exist -- nothing to compile
add_library(minimatrix INTERFACE)

target_include_directories(minimatrix INTERFACE
    ${CMAKE_CURRENT_SOURCE_DIR}/include
)
target_compile_features(minimatrix INTERFACE cxx_std_17)

# Usage:
add_executable(demo demo.cpp)
target_link_libraries(demo PRIVATE minimatrix)  # gets the include path automatically
\end{lstlisting}

\gsub{ALIAS Hedefleri --- Namespace'li İsimler}

Kütüphanelere \code{Ad::Alt} biçiminde takma ad vermek yerleşmiş, profesyonel
bir alışkanlıktır; böylece kütüphaneniz \code{find\_package} ile gelen
hedeflerle aynı görünümü kazanır.

\begin{lstlisting}
add_library(math STATIC src/matrix.cpp)
add_library(MiniMat::math ALIAS math)   # namespaced alias name

# Now it can be used like this (as if it came from find_package):
target_link_libraries(app PRIVATE MiniMat::math)
\end{lstlisting}

\ipucu{\code{Ad::} içeren bir hedef adı, CMake tarafından ``bu bir gerçek
hedeftir'' olarak katı denetlenir: yanlış yazarsanız sessizce \code{-lYanlis}
denemek yerine hemen hata verir. Bu yüzden namespace'li adlar daha güvenlidir.}

% #####################################################################
\gpart{3}{Orta Seviye --- Yapı, Bağımlılık ve Mantık}
% #####################################################################

% =====================================================================
\gsec{Çok Dizinli Proje Düzeni}
% =====================================================================

Gerçek projeler tek bir \code{CMakeLists.txt} dosyasıyla yürütülmez. Her
mantıksal bileşen kendi alt dizininde, kendi \code{CMakeLists.txt} dosyasıyla
yaşar; kök dosya da bunların hepsini \code{add\_subdirectory} ile bir araya
getirir.

\gsub{Önerilen Klasör Yapısı}

\begin{lstlisting}[style=shstyle]
myproject/
|-- CMakeLists.txt          # root: project-wide settings
|-- include/                # public headers (optional centralized location)
|-- src/
|   `-- CMakeLists.txt      # the main application
|-- libs/
|   `-- math/
|       |-- CMakeLists.txt  # the math library
|       |-- include/math/
|       |   `-- matrix.h
|       `-- src/matrix.cpp
|-- tests/
|   `-- CMakeLists.txt      # tests
`-- build/                  # (generated; listed in .gitignore)
\end{lstlisting}

\gsub{Kök CMakeLists.txt}

\begin{lstlisting}[title={\color{ortaGri}\ttfamily CMakeLists.txt (kök)}]
cmake_minimum_required(VERSION 3.20)
project(MyProject VERSION 1.0 LANGUAGES CXX)

# Project-wide settings
set(CMAKE_CXX_STANDARD 17)
set(CMAKE_CXX_STANDARD_REQUIRED ON)

# Collect all executables in one place (optional, but keeps things tidy)
set(CMAKE_RUNTIME_OUTPUT_DIRECTORY ${CMAKE_BINARY_DIR}/bin)

# Emit compile_commands.json (very useful for clangd and IDEs)
set(CMAKE_EXPORT_COMPILE_COMMANDS ON)

# Add subdirectories in order (dependency order matters!)
add_subdirectory(libs/math)   # the library first
add_subdirectory(src)         # then the application that uses it

# Add the tests conditionally
option(BUILD_TESTS "Build the unit tests" ON)
if(BUILD_TESTS)
    enable_testing()
    add_subdirectory(tests)
endif()
\end{lstlisting}

\gsub{Kütüphane Alt Dizini}

\begin{lstlisting}[title={\color{ortaGri}\ttfamily libs/math/CMakeLists.txt}]
add_library(math STATIC
    src/matrix.cpp
    src/vector.cpp
)
add_library(MyProject::math ALIAS math)

target_include_directories(math PUBLIC
    ${CMAKE_CURRENT_SOURCE_DIR}/include
)
target_compile_features(math PUBLIC cxx_std_17)
\end{lstlisting}

\begin{lstlisting}[title={\color{ortaGri}\ttfamily src/CMakeLists.txt}]
add_executable(app main.cpp)
target_link_libraries(app PRIVATE MyProject::math)  # consume the library
\end{lstlisting}

\bilgi{\code{add\_subdirectory} sırası önemlidir: bir hedefe
\code{target\_link\_libraries} ile bağlanmadan önce o hedefin tanımlanmış
olması gerekir. Kütüphaneleri her zaman kullanıcılarından önce ekleyin.}

\miniproj{Yukarıdaki düzeni çalışan hale getirmek}

Eksik kalan tek şey kaynak dosyaları. Aşağıdaki üç dosyayı ekleyince proje
uçtan uca derlenir (testleri şimdilik kapatıyoruz):

\begin{lstlisting}[style=cppstyle,title={\color{ortaGri}\ttfamily libs/math/include/math/matrix.h}]
#pragma once
namespace mp { int add(int a, int b); }
\end{lstlisting}

\begin{lstlisting}[style=cppstyle,title={\color{ortaGri}\ttfamily libs/math/src/matrix.cpp}]
#include "math/matrix.h"
namespace mp { int add(int a, int b) { return a + b; } }
\end{lstlisting}

\begin{lstlisting}[style=cppstyle,title={\color{ortaGri}\ttfamily src/main.cpp}]
#include <iostream>
#include "math/matrix.h"        // include/math yolu MyProject::math'ten geldi
int main() { std::cout << "2 + 3 = " << mp::add(2, 3) << "\n"; }
\end{lstlisting}

\begin{lstlisting}[style=shstyle]
# tests/ klasorunu yazmadigimiz icin BUILD_TESTS=OFF veriyoruz
cmake -S . -B build -G Ninja -D BUILD_TESTS=OFF
cmake --build build
./build/bin/app          # cikti build/bin altina gider (koktaki ayar)
#   2 + 3 = 5
\end{lstlisting}

% =====================================================================
\gsec{Harici Bağımlılıklar: find\_package}
% =====================================================================

Çoğu proje, sistemde kurulu kütüphanelerden (Threads, OpenGL, Boost, SQLite,
fmt gibi) yararlanır. \code{find\_package} bu kütüphaneleri sizin için bulur ve
doğrudan bağlayabileceğiniz hazır hedefler sunar.

\gsub{İki Mod: Config ve Module}

\code{find\_package} iki şekilde arama yapar:

\begin{center}
\begin{tabularx}{\textwidth}{@{}>{\bfseries\color{anaMavi}}l >{\raggedright\arraybackslash}X@{}}
\toprule
\rowcolor{acikMavi}\color{anaMavi}\textbf{Mod} & \color{anaMavi}\textbf{Açıklama}\\
\midrule
Config & Kütüphanenin kendi kurduğu \code{XxxConfig.cmake} dosyasını bulur. Modern kütüphaneler bunu sağlar (en güvenilir).\\
\rowcolor{acikGri} Module & CMake ile gelen veya sizin yazdığınız \code{FindXxx.cmake} bulma betiğini kullanır (eski kütüphaneler için).\\
\bottomrule
\end{tabularx}
\end{center}

\gsub{Örnek: Threads (pthread)}

\begin{lstlisting}
find_package(Threads REQUIRED)     # REQUIRED: stop with an error if not found

add_executable(server server.cpp)
target_link_libraries(server PRIVATE Threads::Threads)
\end{lstlisting}

\gsub{Örnek: fmt Kütüphanesi (Config mod)}

\begin{lstlisting}
find_package(fmt 9.0 REQUIRED)     # require version 9.0 or newer

add_executable(app main.cpp)
target_link_libraries(app PRIVATE fmt::fmt)
\end{lstlisting}

\begin{lstlisting}[style=cppstyle,title={\color{ortaGri}\ttfamily main.cpp}]
#include <fmt/core.h>
int main() {
    // {} are placeholders, filled positionally by the following arguments
    fmt::print("Hello {}! {} + {} = {}\n", "world", 2, 3, 2 + 3);
}
\end{lstlisting}

\gsub{COMPONENTS ve Opsiyonel Kullanım}

\begin{lstlisting}
# Request only specific components of a large package
find_package(Boost 1.75 REQUIRED COMPONENTS filesystem system)
target_link_libraries(app PRIVATE Boost::filesystem Boost::system)

# Without REQUIRED: enable the feature if found, otherwise continue silently
find_package(OpenMP)
if(OpenMP_CXX_FOUND)
    target_link_libraries(app PRIVATE OpenMP::OpenMP_CXX)
    message(STATUS "OpenMP found; parallel build enabled")
else()
    message(STATUS "OpenMP not found; building without parallelism")
endif()
\end{lstlisting}

\notk{CMake'e kütüphanenin nerede olduğunu söylemek için
\code{-DCMAKE\_PREFIX\_PATH=/opt/mylib} kullanın. \code{find\_package} bu
yolları da tarar.}

\gsub{Elle Kurulan Kütüphaneler: CMake Nerede Arar?}

Sık karşılaşılan bir durum: bir kütüphaneyi paket yöneticisinden değil,
kaynaktan derleyip \code{make install} (veya \code{cmake --install}) ile
kurarsınız. Bu genellikle dosyaları \code{/usr/local} altına yerleştirir:

\begin{lstlisting}[style=shstyle]
/usr/local/lib/            # .so / .a  (kutuphane dosyalari)
/usr/local/include/        # .h / .hpp (basliklar)
/usr/local/lib/cmake/Xxx/  # XxxConfig.cmake  <-- find_package BUNU arar
\end{lstlisting}

\code{find\_package(Xxx)}, kütüphane dosyasını (\code{.so}) değil, o kütüphanenin
kurulum sırasında bıraktığı \code{XxxConfig.cmake} dosyasını arar. CMake bu
dosyayı ararken şu \emph{ön ek} (prefix) yollarının her birinin altına bakar:

\begin{center}
\begin{tabularx}{\textwidth}{@{}>{\bfseries\color{anaMavi}}l >{\raggedright\arraybackslash}X@{}}
\toprule
\rowcolor{acikMavi}\color{anaMavi}\textbf{Aranan yol (\code{<prefix>} altında)} & \color{anaMavi}\textbf{Açıklama}\\
\midrule
\code{lib/cmake/Xxx/} & En yaygın konum.\\
\rowcolor{acikGri} \code{lib/Xxx/cmake/} & Bazı kütüphanelerin tercihi.\\
\code{share/Xxx/cmake/} & Mimariden bağımsız kurulumlar.\\
\rowcolor{acikGri} \code{lib64/cmake/Xxx/} & Bazı 64-bit dağıtımlar.\\
\bottomrule
\end{tabularx}
\end{center}

Buradaki \code{<prefix>}, CMake'in varsayılan olarak taradığı köklerdir:
\code{CMAKE\_SYSTEM\_PREFIX\_PATH}. Linux'ta bu liste \code{/usr} ve
\code{/usr/local}'ı \emph{zaten içerir}; bu yüzden \code{/usr/local} altına
düzgün kurulmuş çoğu kütüphane hiçbir ayar yapmadan bulunur.

\gsubsub{Peki neden bazen bulunamıyor?}

\begin{itemize}
\item \textbf{Farklı bir ön eke kuruldu:} Kütüphane \code{/usr/local} yerine
\code{/opt/mylib} ya da \code{\$HOME/.local} gibi standart dışı bir yere
kurulmuştur (\code{-DCMAKE\_INSTALL\_PREFIX} veya \code{./configure --prefix}).
CMake oralara varsayılan olarak bakmaz.
\item \textbf{Config dosyası hiç üretilmemiş:} Kütüphane yalnızca \code{.so} ve
başlık kurar, \code{XxxConfig.cmake} sağlamaz (eski veya Autotools tabanlı
projeler). Bu durumda \code{find\_package}'in Config modu için tutunacak bir
şey yoktur.
\item \textbf{\code{lib} / \code{lib64} uyuşmazlığı:} Dosyalar \code{lib64}
altındayken CMake \code{lib} altına bakıyordur (dağıtıma göre değişir).
\end{itemize}

\gsubsub{Çözüm 1 --- CMAKE\_PREFIX\_PATH (standart dışı ön ek)}

CMake'e ``şu kök(ler)in altına da bak'' demenin en temiz yolu budur.
Verdiğiniz yolun altındaki \code{lib/cmake/...} vb. otomatik taranır:

\begin{lstlisting}[style=shstyle]
# Kutuphane /opt/mylib altina kuruldu diyelim:
cmake -S . -B build -DCMAKE_PREFIX_PATH=/opt/mylib

# Birden fazla yol -- noktali virgulle ayrilir:
cmake -S . -B build -DCMAKE_PREFIX_PATH="/opt/mylib;$HOME/.local"
\end{lstlisting}

Aynı şeyi kalıcı olarak ortam değişkeni ile de verebilirsiniz; CMake bunu okur:

\begin{lstlisting}[style=shstyle]
export CMAKE_PREFIX_PATH=/opt/mylib:$HOME/.local   # (: ile ayirin)
\end{lstlisting}

\gsubsub{Çözüm 2 --- Xxx\_DIR (doğrudan config klasörü)}

Yalnızca \emph{tek} bir paket sorun çıkarıyorsa, o paketin \code{Config.cmake}
dosyasının bulunduğu klasörü doğrudan gösterebilirsiniz. Değişkenin adı
\code{<PaketAdı>\_DIR} biçimindedir:

\begin{lstlisting}[style=shstyle]
cmake -S . -B build \
      -Dfmt_DIR=/opt/fmt/lib/cmake/fmt   # XxxConfig.cmake'in oldugu klasor
\end{lstlisting}

\gsubsub{Çözüm 3 --- Config yoksa: pkg-config'e düşün}

Kütüphane hiç \code{XxxConfig.cmake} sağlamıyorsa ama bir \code{.pc}
(\code{/usr/local/lib/pkgconfig/xxx.pc}) dosyası bırakmışsa, \code{find\_package}
yerine \code{pkg-config} köprüsünü kullanabilirsiniz:

\begin{lstlisting}
find_package(PkgConfig REQUIRED)
pkg_check_modules(XXX REQUIRED IMPORTED_TARGET xxx)  # xxx = .pc dosyasinin adi

add_executable(app main.cpp)
target_link_libraries(app PRIVATE PkgConfig::XXX)    # bayraklar/yollar otomatik
\end{lstlisting}

\gsubsub{Örnek: İki farklı yoldaki iki kütüphane}

Gerçekçi bir senaryo: projeniz iki kütüphaneye bağlı. Biri kaynaktan derleyip
\code{/usr/local} altına kurduğunuz kendi \code{mylib}'iniz, diğeri ise
resmi kurulumla \code{/opt/Qt/6.7.0/gcc\_64} altına gelen \textbf{Qt}. İkisi
\emph{farklı ön eklerde} durduğundan, \code{CMAKE\_PREFIX\_PATH}'e her ikisini
de vermeniz gerekir:

\begin{lstlisting}[style=shstyle]
# Her iki paketin de config dosyalari bulunan koklerini ; ile ver:
cmake -S . -B build -G Ninja \
      -DCMAKE_PREFIX_PATH="/usr/local;/opt/Qt/6.7.0/gcc_64"

#   /usr/local           altinda -> lib/cmake/mylib/mylibConfig.cmake
#   /opt/Qt/6.7.0/gcc_64 altinda -> lib/cmake/Qt6/Qt6Config.cmake
\end{lstlisting}

\begin{lstlisting}[title={\color{ortaGri}\ttfamily CMakeLists.txt}]
cmake_minimum_required(VERSION 3.20)
project(IkiKutuphane LANGUAGES CXX)
set(CMAKE_CXX_STANDARD 17)

# find_package her iki koku de CMAKE_PREFIX_PATH sayesinde tarar:
find_package(mylib REQUIRED)                       # /usr/local altindan
find_package(Qt6 REQUIRED COMPONENTS Widgets)      # /opt/Qt/... altindan

add_executable(app main.cpp)
target_link_libraries(app PRIVATE
    mylib::mylib          # kendi kutuphaneniz
    Qt6::Widgets)         # Qt bileseni
\end{lstlisting}

\notk[İpucu]{Yalnızca Qt sorun çıkarıyorsa, \code{CMAKE\_PREFIX\_PATH} yerine
tek bir paketi hedefleyen \code{-DQt6\_DIR=/opt/Qt/6.7.0/gcc\_64/lib/cmake/Qt6}
da işe yarar. \code{CMAKE\_PREFIX\_PATH} \emph{birden çok} paket için, \code{\_DIR}
\emph{tek} bir paket için pratiktir.}

\uyari{\code{/usr/local/lib}'e kurdunuz ama \emph{çalıştırırken}
\code{error while loading shared libraries: libxxx.so} alıyorsanız bu bir CMake
sorunu değildir: dinamik yükleyici bu klasörü taramıyordur. \code{/etc/ld.so.conf.d/}
altına yolu ekleyip \code{sudo ldconfig} çalıştırın veya
\code{export LD\_LIBRARY\_PATH=/usr/local/lib} verin.}

\ipucu{Paket gerçekten nerede aranıyor, göremiyor musunuz?
\code{cmake --debug-find -S . -B build} komutu \code{find\_package}'in
denediği \emph{tüm} yolları tek tek yazdırır; hangi klasöre bakıp bulamadığını
anında görürsünüz.}

\miniproj{\code{find\_package(Threads)} ile çalışan bir çok-iş-parçacıklı program}

\code{Threads}, hemen her sistemde hazır bulunur; ayrıca bir şey kurmanız
gerekmez. Bu yüzden \code{find\_package}'i denemenin en kolay yolu budur.

\begin{lstlisting}[style=shstyle]
threads-demo/
|-- CMakeLists.txt
`-- main.cpp
\end{lstlisting}

\begin{lstlisting}[style=cppstyle,title={\color{ortaGri}\ttfamily main.cpp}]
#include <iostream>
#include <thread>
#include <vector>

int main() {
    std::vector<std::thread> workers;
    for (int i = 0; i < 3; ++i)
        workers.emplace_back([i]{ std::cout << "Is parcacigi " << i << "\n"; });
    for (auto& t : workers) t.join();   // hepsinin bitmesini bekle
}
\end{lstlisting}

\begin{lstlisting}[title={\color{ortaGri}\ttfamily CMakeLists.txt}]
cmake_minimum_required(VERSION 3.20)
project(ThreadsDemo LANGUAGES CXX)
set(CMAKE_CXX_STANDARD 17)

# Sistemdeki is parcacigi destegini bul (Linux'ta pthread'e cozulur)
find_package(Threads REQUIRED)

add_executable(app main.cpp)
target_link_libraries(app PRIVATE Threads::Threads)  # -pthread'i CMake ekler
\end{lstlisting}

\begin{lstlisting}[style=shstyle]
cmake -S . -B build -G Ninja && cmake --build build
./build/app
#   Is parcacigi 0
#   Is parcacigi 1      # sira karisabilir -- is parcaciklari es zamanli calisir
#   Is parcacigi 2
\end{lstlisting}

% =====================================================================
\gsec{Fonksiyonlar ve Makrolar}
% =====================================================================

Bir CMake fonksiyonu, tıpkı herhangi bir programlama dilindeki fonksiyon gibi,
tekrar eden bir işi \emph{bir kez} tanımlayıp \emph{defalarca} çağırmanızı
sağlar. Aynı beş satırlık hedef kurulumunu on ayrı yere kopyalamak yerine, o
mantığı bir fonksiyonun içine koyar ve her seferinde tek satırla çağırırsınız.
Bu bölümde önce bir fonksiyonun yapısını, ardından \emph{gerçek bir projede}
onu nasıl tanımlayıp kullanacağınızı adım adım ele alacağız.

\gsub{function vs macro}

\begin{center}
\begin{tabularx}{\textwidth}{@{}>{\bfseries\color{anaMavi}}l >{\raggedright\arraybackslash}X@{}}
\toprule
\rowcolor{acikMavi}\color{anaMavi}\textbf{function} & \color{anaMavi}\textbf{macro}\\
\midrule
Kendi değişken kapsamı vardır (yereldir). & Kapsamı yoktur; çağrıldığı yere metin gibi genişler.\\
\rowcolor{acikGri} \code{PARENT\_SCOPE} ile dışarı yazar. & Çağıranın değişkenlerini doğrudan değiştirir.\\
Genelde tercih edilir (daha güvenli). & Sadece değişken sızdırmak gerektiğinde.\\
\bottomrule
\end{tabularx}
\end{center}

\gsub{Bir Fonksiyonun Anatomisi: Parametreler ve ARGN}

Bir fonksiyonu \code{function(<isim> <par1> <par2> ...)} ile tanımlarsınız.
Adlandırılmış parametrelerin yanı sıra CMake, her çağrıda birkaç
\textbf{otomatik değişkeni} de sizin için doldurur: \code{ARGC} (argüman
sayısı), \code{ARGV} (tüm argümanlar, tek bir liste olarak), \code{ARGN}
(adlandırılmış parametrelerden \emph{sonra} gelen ``fazladan'' argümanlar) ve
\code{ARGV0}, \code{ARGV1}, \code{...} (indeksle erişim). Bunlardan \code{ARGN}
özellikle önemlidir: değişken sayıda kaynak dosyası alan hedef kurucular tam da
bunun üzerine inşa edilir.

\begin{lstlisting}
function(demo first second)
    # Named positional parameters:
    message(STATUS "first  = ${first}")   # 1st argument
    message(STATUS "second = ${second}")  # 2nd argument

    # Automatic variables available inside every function:
    message(STATUS "ARGC = ${ARGC}")      # total number of arguments
    message(STATUS "ARGV = ${ARGV}")      # ALL arguments, as a ;-list
    message(STATUS "ARGN = ${ARGN}")      # only the EXTRA args (after first, second)
    message(STATUS "ARGV0 = ${ARGV0}")    # index access -> same as 'first'
endfunction()

demo(one two three four)
# -- first  = one
# -- second = two
# -- ARGC = 4
# -- ARGV = one;two;three;four
# -- ARGN = three;four
# -- ARGV0 = one
\end{lstlisting}

\bilgi{Bir fonksiyon \emph{kendi değişken kapsamına} sahiptir: içinde
\code{set(x ...)} yaptığınızda bu \code{x}, fonksiyon bittiğinde yok olur ve
çağıranın \code{x}'ini kirletmez. Bu yüzden hedef kurmak gibi işler için
fonksiyon, makrodan daha güvenlidir.}

\gsub{Neden Fonksiyon? Önce/Sonra}

Diyelim ki projenizde bir \code{examples/} klasörü var ve içinde onlarca küçük
örnek program bulunuyor. Bunların her biri \emph{aynı} kütüphaneye bağlanmalı,
\emph{aynı} C++ standardını ve \emph{aynı} uyarı bayraklarını kullanmalı,
çıktısı da \emph{aynı} klasöre düşmeli. Fonksiyon kullanmazsanız bu, her örnek
için tekrar tekrar yazılan dört satır demektir:

\begin{lstlisting}[title={\color{ortaGri}\ttfamily examples/CMakeLists.txt --- fonksiyonsuz (tekrar dolu)}]
add_executable(gauss gauss.cpp)
target_link_libraries(gauss PRIVATE MyProject::math project_warnings)
target_compile_features(gauss PRIVATE cxx_std_17)
set_target_properties(gauss PROPERTIES RUNTIME_OUTPUT_DIRECTORY ${CMAKE_BINARY_DIR}/examples)

add_executable(fourier fourier.cpp fft.cpp)
target_link_libraries(fourier PRIVATE MyProject::math project_warnings)
target_compile_features(fourier PRIVATE cxx_std_17)
set_target_properties(fourier PROPERTIES RUNTIME_OUTPUT_DIRECTORY ${CMAKE_BINARY_DIR}/examples)

# ... 10 örnek daha, her biri aynı 4 satır -> bakımı imkansız, hataya açık
\end{lstlisting}

Bu mantığı tek bir \code{add\_example} fonksiyonuna taşırsak, aynı iş her örnek
için \textbf{tek satıra} iner:

\begin{lstlisting}[title={\color{ortaGri}\ttfamily examples/CMakeLists.txt --- fonksiyonlu (temiz)}]
add_example(gauss   gauss.cpp)              # tek kaynak
add_example(fourier fourier.cpp fft.cpp)    # birden çok kaynak (ARGN'e gider)
add_example(newton  newton.cpp)
\end{lstlisting}

\gsub{Gerçek Proje: Fonksiyonu Ayrı Bir Modüle Koymak}

Profesyonel projelerde yardımcı fonksiyonlar kök \code{CMakeLists.txt} içine
değil, \code{cmake/} klasöründeki ayrı bir \code{.cmake} modülüne yazılır; kök
dosya da bu modülü \code{include()} ile yükler. Böylece fonksiyon tüm projede,
yani bütün alt dizinlerde kullanılabilir hale gelir. Aşağıda bunun uçtan uca
tam bir örneği var.

\begin{lstlisting}[style=shstyle,title={\color{ortaGri}\ttfamily Proje düzeni}]
myproject/
|-- CMakeLists.txt            # kök: modülü yükler, alt dizinleri ekler
|-- cmake/
|   `-- AddExample.cmake      # yeniden kullanılabilir yardımcı fonksiyon
|-- libs/math/                # MyProject::math kütüphanesi
`-- examples/
    |-- CMakeLists.txt        # sadece add_example(...) çağrıları
    |-- gauss.cpp
    |-- fourier.cpp
    `-- fft.cpp
\end{lstlisting}

Önce fonksiyonun \emph{tanımı} kendi modül dosyasına yazılır:

\begin{lstlisting}[title={\color{ortaGri}\ttfamily cmake/AddExample.cmake}]
# Reusable helper: build ONE example program, already wired to the math
# library, the shared warning flags, and a common output directory.
include_guard(GLOBAL)   # run this file's body at most once, even if included twice

function(add_example name)
    # name  -> the target (executable) name, e.g. "gauss"
    # ARGN  -> every remaining argument = this example's source files
    add_executable(${name} ${ARGN})

    target_link_libraries(${name} PRIVATE
        MyProject::math       # our own library (namespaced ALIAS target)
        project_warnings      # shared warning policy (an INTERFACE target)
    )
    target_compile_features(${name} PRIVATE cxx_std_17)

    set_target_properties(${name} PROPERTIES
        RUNTIME_OUTPUT_DIRECTORY ${CMAKE_BINARY_DIR}/examples
    )

    # Bonus: also register the example as a smoke test -- running it must
    # succeed (exit code 0). One function call now gives us build + test.
    add_test(NAME example_${name} COMMAND ${name})
endfunction()
\end{lstlisting}

Kök dosya, CMake'e ``kendi modüllerimi \code{cmake/} klasöründe ara'' dedikten
sonra fonksiyonu \code{include} ile yükler. Bu \code{include} çağrısından
\emph{sonra} eklenen her alt dizin, artık bu fonksiyona erişebilir:

\begin{lstlisting}[title={\color{ortaGri}\ttfamily CMakeLists.txt (kök) --- ilgili kısım}]
cmake_minimum_required(VERSION 3.20)
project(MyProject VERSION 1.0 LANGUAGES CXX)

# 1) CMake'in modül arama yoluna kendi cmake/ klasörümüzü ekle
list(APPEND CMAKE_MODULE_PATH ${CMAKE_SOURCE_DIR}/cmake)

# 2) Modülü yükle -> artık add_example() her yerde tanımlı
include(AddExample)          # AddExample.cmake dosyasını çalıştırır

enable_testing()             # add_test() çağrılarının kaydedilmesi için

add_subdirectory(libs/math)  # önce kütüphane (add_example onu kullanıyor)
add_subdirectory(examples)   # sonra örnekler
\end{lstlisting}

Son olarak alt dizin, tüm o karmaşayı unutup sadece \emph{ne} istediğini yazar:

\begin{lstlisting}[title={\color{ortaGri}\ttfamily examples/CMakeLists.txt}]
add_example(gauss   gauss.cpp)              # -> hedef "gauss" + test "example_gauss"
add_example(fourier fourier.cpp fft.cpp)    # fazladan kaynak ARGN ile gelir
add_example(newton  newton.cpp)
\end{lstlisting}

\ornek[Ne oluyor?]{\code{add\_example(fourier fourier.cpp fft.cpp)} satırı
çalıştığında CMake, sanki o dört satırı \code{examples/CMakeLists.txt} içine
elle yazmışsınız gibi davranır: \code{name} = \code{"fourier"}, \code{ARGN} =
\code{"fourier.cpp;fft.cpp"} olur; \code{add\_executable} bu iki kaynağı
derler, kütüphane ve uyarılar bağlanır, çıktı \code{build/examples/} altına
konur ve bir de \code{example\_fourier} testi kaydedilir. Onlarca örnek için
bakılacak tek yer artık \code{AddExample.cmake} dosyasıdır.}

\bilgi{Bir fonksiyon \emph{bir kez} (global olarak) tanımlanır ve
tanımlandığı andan itibaren \textbf{tüm alt dizinlerde} görünür --- değişkenlerin
aksine, alt dizine ``kopyalanması'' gerekmez. Bu yüzden yardımcıları kökte
\code{include} etmek, onları tek merkezden tüm projeye yaymanın yoludur.}

\ipucu{Yeniden kullanılabilir fonksiyonları \code{cmake/} klasöründe ayrı
\code{.cmake} dosyalarında toplayın (\code{AddExample.cmake},
\code{Warnings.cmake}\ldots). \code{include\_guard(GLOBAL)} ile dosya iki kez
\code{include} edilse bile fonksiyon ``already defined'' hatası vermez.}

\gsub{İsimli Argümanlar: cmake\_parse\_arguments}

Yukarıdaki \code{add\_example} basit bir örnekti; çünkü fazladan gelen tüm
argümanları (\code{ARGN}) doğrudan kaynak dosyası sayıyordu. Ama fonksiyonunuzun
\emph{seçenekler} (bayraklar), \emph{tek değerli} ve \emph{çok değerli}
anahtarlar alması gerekiyorsa, \code{ARGN}'i elle ayrıştırmak epey zahmetlidir.
CMake'in yerleşik \code{cmake\_parse\_arguments} komutu bu işi sizin yerinize
yapar: her anahtarı \code{<önek>\_<ANAHTAR>} biçiminde bir değişkene yerleştirir.

\begin{lstlisting}
function(add_app name)
    # cmake_parse_arguments(<prefix> <options> <one_value> <multi_value> ${ARGN})
    #   options    "GUI"          -> ARG_GUI       TRUE/FALSE (bayrak var mı?)
    #   one_value  "VERSION"      -> ARG_VERSION   tek bir değer
    #   multi      "SOURCES;LIBS" -> ARG_SOURCES / ARG_LIBS   birer liste
    cmake_parse_arguments(ARG "GUI" "VERSION" "SOURCES;LIBS" ${ARGN})

    add_executable(${name} ${ARG_SOURCES})
    if(ARG_LIBS)
        target_link_libraries(${name} PRIVATE ${ARG_LIBS})
    endif()
    if(ARG_VERSION)
        target_compile_definitions(${name} PRIVATE APP_VERSION="${ARG_VERSION}")
    endif()
    if(ARG_GUI)                       # sadece GUI anahtarı verildiyse tanımla
        target_compile_definitions(${name} PRIVATE USE_GUI)
    endif()
endfunction()

# Çağrı: anahtarların sırası önemli değildir, okunması kolaydır
add_app(app
    GUI                              # bir bayrak (değer almaz)
    VERSION 1.4                      # tek değerli anahtar
    SOURCES main.cpp ui.cpp          # çok değerli anahtar
    LIBS    math Threads::Threads    # çok değerli anahtar
)
\end{lstlisting}

\notk{Böyle bir imza, çağrı yerinde kodun \emph{niyetini} okunur kılar:
\code{add\_app(app GUI VERSION 1.4 SOURCES ... LIBS ...)} satırına bakan biri,
argümanların ne anlama geldiğini sırayı ezberlemeden anlar. Büyük projelerde
kendi hedef kurucularınızı hep bu kalıpla yazmak, \code{CMakeLists.txt}
dosyalarını okunur tutar.}

\miniproj{Küçük ama çalışan bir \code{add\_example} fonksiyonu}

Yukarıdaki örnek büyük bir projeyi anlatıyordu. İşte fikri kanıtlayan, hemen
derlenebilen en küçük hali: tek fonksiyon, iki örnek program.

\begin{lstlisting}[style=shstyle]
funcs-demo/
|-- CMakeLists.txt
|-- hello.cpp
`-- bye.cpp
\end{lstlisting}

\begin{lstlisting}[style=cppstyle,title={\color{ortaGri}\ttfamily hello.cpp / bye.cpp}]
// hello.cpp
#include <iostream>
int main() { std::cout << "Merhaba fonksiyon!\n"; }

// bye.cpp
#include <iostream>
int main() { std::cout << "Gule gule!\n"; }
\end{lstlisting}

\begin{lstlisting}[title={\color{ortaGri}\ttfamily CMakeLists.txt}]
cmake_minimum_required(VERSION 3.20)
project(FuncsDemo LANGUAGES CXX)

# Tekrar eden hedef kurulumunu tek bir fonksiyona sariyoruz
function(add_example name)
    add_executable(${name} ${ARGN})     # ARGN = fazladan gelen kaynak dosyalar
    target_compile_features(${name} PRIVATE cxx_std_17)
    set_target_properties(${name} PROPERTIES
        RUNTIME_OUTPUT_DIRECTORY ${CMAKE_BINARY_DIR}/examples)
endfunction()

# Artik her ornek tek satir:
add_example(hello hello.cpp)
add_example(bye   bye.cpp)
\end{lstlisting}

\begin{lstlisting}[style=shstyle]
cmake -S . -B build -G Ninja && cmake --build build
./build/examples/hello     #   Merhaba fonksiyon!
./build/examples/bye       #   Gule gule!
\end{lstlisting}

% =====================================================================
\gsec{Generator Expressions}
% =====================================================================

Generator expression (kısaca ``genex''), \code{\$<...>} biçiminde yazılan ve
\emph{yapılandırma anında değil, üretim anında} değerlendirilen ifadelerdir.
Derleyiciye, derleme tipine (Debug/Release) ya da hedefin özelliklerine göre
farklı değerler üretmenin tek doğru yolu bunlardır.

\gsub{Neden Gerekir?}

Bir \code{if(CMAKE\_BUILD\_TYPE STREQUAL "Debug")} kontrolü, çok yapılandırmalı
generator'larda (Visual Studio, Xcode) işe yaramaz; çünkü oralarda derleme tipi
\emph{derleme anında} seçilir, yapılandırma anında henüz belli değildir. Genex
tam da bu sorunu çözer.

\gsub{Derleyiciye Göre Bayrak}

\begin{lstlisting}
add_executable(app main.cpp)

target_compile_options(app PRIVATE
    # these flags only when the compiler is GCC or Clang
    $<$<CXX_COMPILER_ID:GNU,Clang>:-Wall -Wextra -Wpedantic>
    # this flag only when the compiler is MSVC
    $<$<CXX_COMPILER_ID:MSVC>:/W4>
)
\end{lstlisting}

\gsub{Derleme Tipine Göre Tanım}

\begin{lstlisting}
target_compile_definitions(app PRIVATE
    # defined only in Debug builds
    $<$<CONFIG:Debug>:DEBUG_MODE ENABLE_ASSERTS>
    # defined only in Release builds
    $<$<CONFIG:Release>:NDEBUG FAST_MATH>
)
\end{lstlisting}

\gsub{Sık Kullanılan Genex Kalıpları}

\begin{center}
\begin{tabularx}{\textwidth}{@{}>{\ttfamily\footnotesize\color{koyuGri}}l >{\raggedright\arraybackslash}X@{}}
\toprule
\rowcolor{acikMavi}\color{anaMavi}\textbf{İfade} & \color{anaMavi}\textbf{Anlamı}\\
\midrule
\$<CONFIG:Debug> & Derleme tipi Debug ise 1, değilse 0.\\
\rowcolor{acikGri} \$<CXX\_COMPILER\_ID:GNU> & Derleyici GCC ise 1.\\
\$<PLATFORM\_ID:Windows> & Hedef platform Windows ise 1.\\
\rowcolor{acikGri} \$<BUILD\_INTERFACE:...> & Sadece derleme ağacında (install değil).\\
\$<INSTALL\_INTERFACE:...> & Sadece kurulumdan sonra.\\
\rowcolor{acikGri} \$<TARGET\_FILE:app> & Hedefin üretilen dosya yolu.\\
\bottomrule
\end{tabularx}
\end{center}

\gsub{BUILD\_INTERFACE / INSTALL\_INTERFACE}

Bir kütüphane kurulduğunda include yolu değişir: kaynak ağacındaki konumdan
kurulum dizinine kayar ($\to$). Bu ikili genex, her iki durumu da doğru şekilde
ele alır; kendi paketinizi yazarken (Kısım 4) bu ayrım kritiktir.

\begin{lstlisting}
target_include_directories(math PUBLIC
    $<BUILD_INTERFACE:${CMAKE_CURRENT_SOURCE_DIR}/include>  # while building here
    $<INSTALL_INTERFACE:include>                            # after installation
)
\end{lstlisting}

\miniproj{Genex'in etkisini gözle görmek: Debug ile Release farklı davransın}

Aşağıdaki program, derleme tipine göre \emph{farklı} bir mesaj yazar --- çünkü
\code{DEBUG\_MODE} makrosunu yalnızca Debug yapıya genex ile ekliyoruz. Kaynak
kod hiç değişmeden, çıktı yapıya göre değişir.

\begin{lstlisting}[style=cppstyle,title={\color{ortaGri}\ttfamily main.cpp}]
#include <iostream>
int main() {
#ifdef DEBUG_MODE
    std::cout << "Debug yapisi: ayrintili gunluk acik\n";
#else
    std::cout << "Release yapisi: hizli mod\n";
#endif
}
\end{lstlisting}

\begin{lstlisting}[title={\color{ortaGri}\ttfamily CMakeLists.txt}]
cmake_minimum_required(VERSION 3.20)
project(GenexDemo LANGUAGES CXX)

add_executable(app main.cpp)

# DEBUG_MODE makrosu SADECE Debug yapida tanimlanir
target_compile_definitions(app PRIVATE $<$<CONFIG:Debug>:DEBUG_MODE>)

# Uyari bayraklari SADECE GCC/Clang'te eklenir (MSVC bunlari reddederdi)
target_compile_options(app PRIVATE
    $<$<CXX_COMPILER_ID:GNU,Clang>:-Wall;-Wextra>)
\end{lstlisting}

\begin{lstlisting}[style=shstyle]
# Ayni kaynak, iki farkli yapi -> iki farkli cikti
cmake -S . -B build-dbg -G Ninja -D CMAKE_BUILD_TYPE=Debug
cmake --build build-dbg && ./build-dbg/app
#   Debug yapisi: ayrintili gunluk acik

cmake -S . -B build-rel -G Ninja -D CMAKE_BUILD_TYPE=Release
cmake --build build-rel && ./build-rel/app
#   Release yapisi: hizli mod
\end{lstlisting}

% =====================================================================
\gsec{Derleme Tipleri, Uyarılar ve Sanitizer'lar}
% =====================================================================

\gsub{CMAKE\_BUILD\_TYPE}

Tek yapılandırmalı generator'larda (Make, Ninja) derleme tipini
\code{CMAKE\_BUILD\_TYPE} değişkeni belirler. Her tip, kendine özgü optimizasyon
ve hata ayıklama bayraklarıyla gelir.

\begin{center}
\begin{tabularx}{\textwidth}{@{}>{\bfseries\color{anaMavi}}l >{\raggedright\arraybackslash}X >{\ttfamily\footnotesize}X@{}}
\toprule
\rowcolor{acikMavi}\color{anaMavi}\textbf{Tip} & \color{anaMavi}\textbf{Amaç} & \color{anaMavi}\textbf{Tipik bayraklar}\\
\midrule
Debug & Hata ayıklama; optimizasyon yok, sembol var. & -g -O0\\
\rowcolor{acikGri} Release & Üretim; tam optimizasyon. & -O3 -DNDEBUG\\
RelWithDebInfo & Optimize + sembollü (profil/crash için). & -O2 -g -DNDEBUG\\
\rowcolor{acikGri} MinSizeRel & En küçük ikili. & -Os -DNDEBUG\\
\bottomrule
\end{tabularx}
\end{center}

\begin{lstlisting}[style=shstyle]
cmake -S . -B build -G Ninja -D CMAKE_BUILD_TYPE=Release
cmake --build build
\end{lstlisting}

\ipucu{Kullanıcı bir tip belirtmezse boş kalır (optimizasyonsuz). Kök
dosyanıza makul bir varsayılan koymak iyidir:}

\begin{lstlisting}
if(NOT CMAKE_BUILD_TYPE AND NOT CMAKE_CONFIGURATION_TYPES)
    set(CMAKE_BUILD_TYPE "RelWithDebInfo" CACHE STRING "" FORCE)
endif()
\end{lstlisting}

\gsub{Uyarıları Merkezi Açmak}

Uyarıları tüm projeye dağıtmak yerine, bir INTERFACE hedefinde toplayıp yalnızca
istediğiniz hedeflere bağlamak modern bir yaklaşımdır.

\begin{lstlisting}
# A virtual (INTERFACE) target that carries the "warning policy"
add_library(project_warnings INTERFACE)
target_compile_options(project_warnings INTERFACE
    $<$<CXX_COMPILER_ID:GNU,Clang>:-Wall;-Wextra;-Wpedantic;-Wshadow>
    $<$<CXX_COMPILER_ID:MSVC>:/W4>
)

# Apply it to any target you like:
target_link_libraries(app PRIVATE project_warnings)
\end{lstlisting}

\gsub{Sanitizer'lar (ASan, UBSan)}

Bellek hatalarını ve tanımsız davranışları yakalamakta sanitizer'ların değeri
ölçülemez. Bunları bir seçeneğin arkasına alalım.

\begin{lstlisting}
option(ENABLE_ASAN "Address + UB sanitizer" OFF)

if(ENABLE_ASAN)
    add_compile_options(-fsanitize=address,undefined -fno-omit-frame-pointer)
    add_link_options(-fsanitize=address,undefined)
endif()
\end{lstlisting}

\begin{lstlisting}[style=shstyle]
cmake -S . -B build -G Ninja -D ENABLE_ASAN=ON -D CMAKE_BUILD_TYPE=Debug
cmake --build build
./build/app     # prints a detailed report if a memory error occurs
\end{lstlisting}

\miniproj{Sanitizer bir hatayı yakalarken görmek}

Kasten hatalı bir program yazalım: dizinin sınırının bir dışına yazıyoruz.
Normalde bu hata sessizce geçebilir; ASan ise programı durdurup tam yerini
gösterir.

\begin{lstlisting}[style=cppstyle,title={\color{ortaGri}\ttfamily main.cpp}]
int main() {
    int a[3] = {1, 2, 3};
    a[3] = 42;        // HATA: gecerli indeksler 0..2, a[3] siniri asiyor
    return a[0];
}
\end{lstlisting}

\begin{lstlisting}[title={\color{ortaGri}\ttfamily CMakeLists.txt}]
cmake_minimum_required(VERSION 3.20)
project(AsanDemo LANGUAGES CXX)

option(ENABLE_ASAN "Address + UB sanitizer" OFF)

add_executable(app main.cpp)
if(ENABLE_ASAN)
    target_compile_options(app PRIVATE -fsanitize=address,undefined -g)
    target_link_options(app    PRIVATE -fsanitize=address,undefined)
endif()
\end{lstlisting}

\begin{lstlisting}[style=shstyle]
cmake -S . -B build -G Ninja -D ENABLE_ASAN=ON -D CMAKE_BUILD_TYPE=Debug
cmake --build build
./build/app
#   main.cpp:3:10: runtime error: index 3 out of bounds for type 'int [3]'
#   ==12345==ERROR: AddressSanitizer: stack-buffer-overflow ...
#       WRITE of size 4 ...  #0 in main main.cpp:3     <- hatanin tam satiri
\end{lstlisting}

\ipucu{Sanitizer'lar programı biraz yavaşlatır, o yüzden yalnızca geliştirme ve
test yaparken (\code{-D ENABLE\_ASAN=ON}) açın; yayımladığınız Release yapıda
kapalı kalsın.}

% #####################################################################
\gpart{4}{İleri Seviye --- Bağımlılık Çekme, Test, Paketleme}
% #####################################################################

% =====================================================================
\gsec{FetchContent: Bağımlılığı Projeye Çekmek}
% =====================================================================

\code{find\_package}, sistemde zaten kurulu olan kütüphaneleri bulur. Peki ya
kullanıcının kütüphaneyi önceden kurmuş olmasını istemiyorsanız?
\code{FetchContent}, bir bağımlılığı derleme sırasında indirir ve doğrudan
projenize katar; kullanıcı yalnızca \code{cmake} çalıştırır, gerisini CMake
halleder.

\begin{lstlisting}
include(FetchContent)

# Declare the dependency (git repository + a fixed tag)
FetchContent_Declare(
    fmt
    GIT_REPOSITORY https://github.com/fmtlib/fmt.git
    GIT_TAG        10.2.1        # always pin a fixed version/tag!
)

# Download, configure and make the dependency's targets available
FetchContent_MakeAvailable(fmt)

add_executable(app main.cpp)
target_link_libraries(app PRIVATE fmt::fmt)   # just like a find_package target
\end{lstlisting}

\uyari{\code{GIT\_TAG} olarak \code{main}/\code{master} gibi hareketli dallar
\emph{kullanmayın}. Her zaman sabit bir sürüm etiketi (\code{10.2.1}) veya
commit hash'i verin --- aksi halde derleme tekrarlanabilir olmaz ve bir gün
beklenmedik şekilde kırılır.}

\gsub{FetchContent + find\_package Hibrit (3.24+)}

CMake 3.24 ile birlikte bir bağımlılık önce sistemde \code{find\_package} ile
aranabilir, bulunamazsa indirilebilir. Böylece iki yaklaşımın da en iyi yanını
bir arada kullanmış olursunuz.

\begin{lstlisting}
FetchContent_Declare(
    fmt
    GIT_REPOSITORY https://github.com/fmtlib/fmt.git
    GIT_TAG        10.2.1
    FIND_PACKAGE_ARGS 10  # try find_package(fmt 10) first; download if missing
)
FetchContent_MakeAvailable(fmt)
\end{lstlisting}

\notk{\code{FetchContent} bağımlılığı sizinle birlikte \emph{kaynak koddan}
derler --- ilk derleme yavaş olabilir ama tam kontrol ve sürüm garantisi verir.
Büyük bağımlılıklar (ör. Qt) için sistem paketi + \code{find\_package} genelde
daha pratiktir.}

\miniproj{fmt kütüphanesini indirip kullanan tam proje}

Bu projeyi kurmak için sistemde \code{fmt} kurulu olması gerekmez; CMake ilk
derlemede onu GitHub'dan çeker (bu adım için internet gerekir).

\begin{lstlisting}[style=shstyle]
fetch-demo/
|-- CMakeLists.txt
`-- main.cpp
\end{lstlisting}

\begin{lstlisting}[style=cppstyle,title={\color{ortaGri}\ttfamily main.cpp}]
#include <fmt/core.h>
int main() {
    // {} yer tutuculari, sonraki argumanlarla sirayla doldurulur
    fmt::print("Merhaba {}! {} + {} = {}\n", "fmt", 2, 3, 2 + 3);
}
\end{lstlisting}

\begin{lstlisting}[title={\color{ortaGri}\ttfamily CMakeLists.txt}]
cmake_minimum_required(VERSION 3.20)
project(FetchDemo LANGUAGES CXX)
set(CMAKE_CXX_STANDARD 17)

include(FetchContent)
FetchContent_Declare(fmt
    GIT_REPOSITORY https://github.com/fmtlib/fmt.git
    GIT_TAG        10.2.1)          # her zaman sabit surum etiketi ver
FetchContent_MakeAvailable(fmt)     # indir, yapilandir, hedeflerini hazir et

add_executable(app main.cpp)
target_link_libraries(app PRIVATE fmt::fmt)
\end{lstlisting}

\begin{lstlisting}[style=shstyle]
cmake -S . -B build -G Ninja    # fmt bu asamada indirilir (ilk sefer yavas)
cmake --build build
./build/app
#   Merhaba fmt! 2 + 3 = 5
\end{lstlisting}

% =====================================================================
\gsec{Test: CTest ile}
% =====================================================================

CMake, kendi test koşucusu \code{ctest} ile birlikte gelir. Testlerinizi
kaydeder, çalıştırır ve sonuçlarını raporlar.

\gsub{Temel Test Kaydı}

\begin{lstlisting}[title={\color{ortaGri}\ttfamily kök CMakeLists.txt}]
enable_testing()             # enable testing (call once, in the root list file)
add_subdirectory(tests)
\end{lstlisting}

\begin{lstlisting}[title={\color{ortaGri}\ttfamily tests/CMakeLists.txt}]
add_executable(math_test math_test.cpp)
target_link_libraries(math_test PRIVATE MyProject::math)

# Register a test: a name + the command to run
add_test(NAME MatrixMultiply COMMAND math_test)

# Decide pass/fail from the program's output instead of its exit code
add_test(NAME VersionCheck COMMAND app --version)
set_tests_properties(VersionCheck PROPERTIES
    PASS_REGULAR_EXPRESSION "1\\.0\\.0")
\end{lstlisting}

\begin{lstlisting}[style=cppstyle,title={\color{ortaGri}\ttfamily math\_test.cpp (basit iddia)}]
#include <cassert>
#include "math/matrix.h"

int main() {
    Matrix a = Matrix::identity(3);
    assert(a(0, 0) == 1.0 && a(0, 1) == 0.0);
    // If an assert fails it aborts the program with a non-zero exit code,
    // and CTest reports the test as "failed".
    return 0;
}
\end{lstlisting}

\begin{lstlisting}[style=shstyle]
cmake --build build
cd build && ctest --output-on-failure   # run all tests, show output on failure
ctest -R Matrix                          # only tests whose name matches "Matrix"
ctest -j 8                               # run 8 tests in parallel
ctest --rerun-failed --output-on-failure # re-run only the tests that failed
\end{lstlisting}

\gsub{GoogleTest ile Entegrasyon}

\begin{lstlisting}
include(FetchContent)
FetchContent_Declare(googletest
    GIT_REPOSITORY https://github.com/google/googletest.git
    GIT_TAG        v1.14.0)
FetchContent_MakeAvailable(googletest)

enable_testing()
add_executable(all_tests test_matrix.cpp test_vector.cpp)
target_link_libraries(all_tests PRIVATE MyProject::math GTest::gtest_main)

# Discover each GoogleTest TEST() case as a separate CTest test
include(GoogleTest)
gtest_discover_tests(all_tests)
\end{lstlisting}

\miniproj{Kendi kendine yeten, çalışan bir test projesi}

Harici bir kütüphane olmadan, sadece \code{assert} kullanan küçük bir test.
Test başarısız olursa program sıfırdan farklı bir kodla çıkar ve CTest onu
``failed'' olarak işaretler.

\begin{lstlisting}[style=shstyle]
ctest-demo/
|-- CMakeLists.txt
`-- test_math.cpp
\end{lstlisting}

\begin{lstlisting}[style=cppstyle,title={\color{ortaGri}\ttfamily test\_math.cpp}]
#include <cassert>
static int add(int a, int b) { return a + b; }

int main() {
    assert(add(2, 3) == 5);
    assert(add(-1, 1) == 0);
    return 0;   // tum assert'ler gecerse 0 = basari
}
\end{lstlisting}

\begin{lstlisting}[title={\color{ortaGri}\ttfamily CMakeLists.txt}]
cmake_minimum_required(VERSION 3.20)
project(CTestDemo LANGUAGES CXX)

enable_testing()                       # testleri etkinlestir (kokte bir kez)
add_executable(math_test test_math.cpp)
add_test(NAME AddWorks COMMAND math_test)   # bir test kaydet: ad + calistirilacak komut
\end{lstlisting}

\begin{lstlisting}[style=shstyle]
cmake -S . -B build -G Ninja && cmake --build build
ctest --test-dir build --output-on-failure
#   Test project .../build
#       Start 1: AddWorks
#   1/1 Test #1: AddWorks .........................   Passed    0.00 sec
#   100% tests passed, 0 tests failed out of 1
\end{lstlisting}

% =====================================================================
\gsec{Kurulum (install) ve Paketleme (CPack)}
% =====================================================================

\gsub{install Komutu}

\code{install()}, \code{cmake -{}-install} çalıştırıldığında hangi dosyaların
nereye kopyalanacağını belirler.

\begin{lstlisting}
# Install the executables and libraries
install(TARGETS app math
    RUNTIME DESTINATION bin        # executables -> <prefix>/bin
    LIBRARY DESTINATION lib        # shared libs (.so/.dll) -> <prefix>/lib
    ARCHIVE DESTINATION lib        # static libs (.a/.lib)  -> <prefix>/lib
)

# Install the public headers
install(DIRECTORY libs/math/include/
    DESTINATION include            # -> <prefix>/include
    FILES_MATCHING PATTERN "*.h"   # copy only header files, keep the tree
)
\end{lstlisting}

\begin{lstlisting}[style=shstyle]
cmake -S . -B build -D CMAKE_INSTALL_PREFIX=/opt/myproject
cmake --build build
cmake --install build              # copies the files into /opt/myproject
\end{lstlisting}

\gsub{CPack ile Dağıtım Paketi}

CPack, projenizi \code{.tar.gz}, \code{.deb}, \code{.rpm}, Windows kurulumu
gibi paketlere dönüştürür.

\begin{lstlisting}
# Place this after your install() rules:
set(CPACK_PACKAGE_NAME "MyProject")
set(CPACK_PACKAGE_VERSION "1.0.0")
set(CPACK_PACKAGE_CONTACT "tolga@example.com")
set(CPACK_GENERATOR "TGZ;DEB")     # which package formats to produce
include(CPack)                     # must be included last
\end{lstlisting}

\begin{lstlisting}[style=shstyle]
cmake --build build
cd build && cpack                  # produces MyProject-1.0.0-Linux.tar.gz, .deb, ...
\end{lstlisting}

\miniproj{Küçük bir programı gerçekten kurmak}

Bir programı bir hedef dizine kurup, kurulan dosya ağacının nasıl göründüğünü
görelim. (Sistem dizinlerini kirletmemek için hedef olarak proje içinde bir
\code{stage/} klasörü seçiyoruz.)

\begin{lstlisting}[style=cppstyle,title={\color{ortaGri}\ttfamily main.cpp}]
#include <iostream>
int main() { std::cout << "kuruldum!\n"; }
\end{lstlisting}

\begin{lstlisting}[title={\color{ortaGri}\ttfamily CMakeLists.txt}]
cmake_minimum_required(VERSION 3.20)
project(InstallDemo LANGUAGES CXX)

add_executable(hello main.cpp)

# hello ikilisi <prefix>/bin altina kurulsun
install(TARGETS hello RUNTIME DESTINATION bin)
\end{lstlisting}

\begin{lstlisting}[style=shstyle]
cmake -S . -B build -G Ninja
cmake --build build
cmake --install build --prefix ./stage    # dosyalari ./stage altina kopyalar

./stage/bin/hello       #   kuruldum!
find stage              #   stage
                        #   stage/bin
                        #   stage/bin/hello
\end{lstlisting}

% =====================================================================
\gsec{Kendi Config Paketini Yazmak}
% =====================================================================

Başkalarının kütüphanenizi \code{find\_package(MyProject)} ile kullanabilmesi
için bir ``config paketi'' hazırlayıp kurmanız gerekir. Bu adım, profesyonel
kütüphane dağıtımının anahtarıdır.

\begin{lstlisting}
include(GNUInstallDirs)
include(CMakePackageConfigHelpers)

# 1) Install the target and record it in an export set (EXPORT)
install(TARGETS math
    EXPORT MyProjectTargets
    ARCHIVE DESTINATION ${CMAKE_INSTALL_LIBDIR}
    LIBRARY DESTINATION ${CMAKE_INSTALL_LIBDIR}
    INCLUDES DESTINATION ${CMAKE_INSTALL_INCLUDEDIR}
)
install(DIRECTORY libs/math/include/ DESTINATION ${CMAKE_INSTALL_INCLUDEDIR})

# 2) Install the export set as a .cmake file (with a namespace prefix)
install(EXPORT MyProjectTargets
    FILE MyProjectTargets.cmake
    NAMESPACE MyProject::
    DESTINATION ${CMAKE_INSTALL_LIBDIR}/cmake/MyProject
)

# 3) Generate the version file (used for version checks in find_package)
write_basic_package_version_file(
    "${CMAKE_CURRENT_BINARY_DIR}/MyProjectConfigVersion.cmake"
    VERSION ${PROJECT_VERSION}
    COMPATIBILITY SameMajorVersion
)

# 4) Generate the config file from a template and install it
configure_package_config_file(
    "${CMAKE_CURRENT_SOURCE_DIR}/cmake/MyProjectConfig.cmake.in"
    "${CMAKE_CURRENT_BINARY_DIR}/MyProjectConfig.cmake"
    INSTALL_DESTINATION ${CMAKE_INSTALL_LIBDIR}/cmake/MyProject
)
install(FILES
    "${CMAKE_CURRENT_BINARY_DIR}/MyProjectConfig.cmake"
    "${CMAKE_CURRENT_BINARY_DIR}/MyProjectConfigVersion.cmake"
    DESTINATION ${CMAKE_INSTALL_LIBDIR}/cmake/MyProject
)
\end{lstlisting}

\begin{lstlisting}[title={\color{ortaGri}\ttfamily cmake/MyProjectConfig.cmake.in}]
@PACKAGE_INIT@
include("${CMAKE_CURRENT_LIST_DIR}/MyProjectTargets.cmake")
\end{lstlisting}

Artık başka bir proje şöyle kullanabilir:

\begin{lstlisting}
find_package(MyProject 1.0 REQUIRED)
target_link_libraries(other_app PRIVATE MyProject::math)
\end{lstlisting}

% =====================================================================
\gsec{CMakePresets.json --- Yapılandırmayı Sabitlemek}
% =====================================================================

Uzun \code{-D...} komutlarını her seferinde yeniden yazmak yerine,
\code{CMakePresets.json} dosyasında hazır yapılandırmalar tanımlarsınız. IDE'ler
de bu dosyayı okuyup aynı ayarları kullanır.

\begin{lstlisting}[style=cppstyle,title={\color{ortaGri}\ttfamily CMakePresets.json},morekeywords={}]
{
  "version": 3,
  "configurePresets": [
    {
      "name": "debug",
      "displayName": "Debug (Ninja)",
      "generator": "Ninja",
      "binaryDir": "${sourceDir}/build/debug",
      "cacheVariables": {
        "CMAKE_BUILD_TYPE": "Debug",
        "ENABLE_ASAN": "ON"
      }
    },
    {
      "name": "release",
      "generator": "Ninja",
      "binaryDir": "${sourceDir}/build/release",
      "cacheVariables": { "CMAKE_BUILD_TYPE": "Release" }
    }
  ],
  "buildPresets": [
    { "name": "debug",   "configurePreset": "debug" },
    { "name": "release", "configurePreset": "release" }
  ]
}
\end{lstlisting}

\begin{lstlisting}[style=shstyle]
cmake --preset debug        # configure with all the preset's settings
cmake --build --preset debug
cmake --list-presets        # list the available presets
\end{lstlisting}

\ipucu{\code{CMakePresets.json}'u git'e ekleyin: ekipteki herkes aynı
yapılandırmayı tek komutla kullanır. Kişisel ayarlar için
\code{CMakeUserPresets.json} (git'e eklenmez) kullanılır.}

\miniproj{Tek komutla yapılandırıp derlenen bir preset projesi}

\begin{lstlisting}[style=shstyle]
preset-demo/
|-- CMakeLists.txt
|-- CMakePresets.json
`-- main.cpp
\end{lstlisting}

\begin{lstlisting}[style=cppstyle,title={\color{ortaGri}\ttfamily main.cpp}]
#include <iostream>
int main() { std::cout << "merhaba preset\n"; }
\end{lstlisting}

\begin{lstlisting}[title={\color{ortaGri}\ttfamily CMakeLists.txt}]
cmake_minimum_required(VERSION 3.20)
project(PresetDemo LANGUAGES CXX)
add_executable(app main.cpp)
\end{lstlisting}

\begin{lstlisting}[style=cppstyle,title={\color{ortaGri}\ttfamily CMakePresets.json},morekeywords={}]
{
  "version": 3,
  "configurePresets": [
    { "name": "debug", "generator": "Ninja",
      "binaryDir": "${sourceDir}/build/debug",
      "cacheVariables": { "CMAKE_BUILD_TYPE": "Debug" } }
  ],
  "buildPresets": [ { "name": "debug", "configurePreset": "debug" } ]
}
\end{lstlisting}

\begin{lstlisting}[style=shstyle]
cmake --preset debug         # uzun -D... satiri yok; hepsi preset'ten gelir
cmake --build --preset debug
./build/debug/app
#   merhaba preset
\end{lstlisting}

% =====================================================================
\gsec{Toolchain Dosyaları ve Çapraz Derleme}
% =====================================================================

Başka bir mimari için derleme --- örneğin Linux'ta çalışırken ARM için çıktı
üretme --- bir ``toolchain dosyası'' aracılığıyla yapılır. Bu dosya, hedef
sistemin derleyicisini ve gerekli ayarlarını tanımlar.

\begin{lstlisting}[title={\color{ortaGri}\ttfamily arm-toolchain.cmake}]
# Describe the target system (this also enables cross-compilation mode)
set(CMAKE_SYSTEM_NAME Linux)
set(CMAKE_SYSTEM_PROCESSOR arm)

# The cross compilers
set(CMAKE_C_COMPILER   arm-linux-gnueabihf-gcc)
set(CMAKE_CXX_COMPILER arm-linux-gnueabihf-g++)

# The target root directory (sysroot)
set(CMAKE_FIND_ROOT_PATH /usr/arm-linux-gnueabihf)

# Search programs on the host, but libraries/headers only in the target sysroot
set(CMAKE_FIND_ROOT_PATH_MODE_PROGRAM NEVER)
set(CMAKE_FIND_ROOT_PATH_MODE_LIBRARY ONLY)
set(CMAKE_FIND_ROOT_PATH_MODE_INCLUDE ONLY)
\end{lstlisting}

\begin{lstlisting}[style=shstyle]
cmake -S . -B build-arm -G Ninja \
      -D CMAKE_TOOLCHAIN_FILE=arm-toolchain.cmake
cmake --build build-arm
\end{lstlisting}

% #####################################################################
\gpart{5}{Pratik --- Tam Projeler ve Başvuru}
% #####################################################################

% =====================================================================
\gsec{Tam Örnek: Çok Modüllü C++ Projesi}
% =====================================================================

Şimdiye kadar öğrendiğimiz her şeyi çalışan tek bir projede bir araya
getirelim: bir \code{math} kütüphanesi, onu kullanan bir \code{app} ve testler.

\gsub{Dosya Yapısı}

\begin{lstlisting}[style=shstyle]
calc/
|-- CMakeLists.txt
|-- CMakePresets.json
|-- libs/math/
|   |-- CMakeLists.txt
|   |-- include/math/calculator.h
|   `-- src/calculator.cpp
|-- app/
|   |-- CMakeLists.txt
|   `-- main.cpp
`-- tests/
    |-- CMakeLists.txt
    `-- calc_test.cpp
\end{lstlisting}

\gsub{Kütüphane Kaynağı}

\begin{lstlisting}[style=cppstyle,title={\color{ortaGri}\ttfamily libs/math/include/math/calculator.h}]
#pragma once
namespace calc {
    double add(double a, double b);
    double multiply(double a, double b);
    long factorial(int n);
}
\end{lstlisting}

\begin{lstlisting}[style=cppstyle,title={\color{ortaGri}\ttfamily libs/math/src/calculator.cpp}]
#include "math/calculator.h"

namespace calc {
    double add(double a, double b)      { return a + b; }
    double multiply(double a, double b) { return a * b; }
    long factorial(int n) {
        long r = 1;
        for (int i = 2; i <= n; ++i) r *= i;
        return r;
    }
}
\end{lstlisting}

\gsub{Tüm CMakeLists.txt Dosyaları}

\begin{lstlisting}[title={\color{ortaGri}\ttfamily CMakeLists.txt (kök)}]
cmake_minimum_required(VERSION 3.21)
project(Calc VERSION 1.0 LANGUAGES CXX)

set(CMAKE_CXX_STANDARD 17)
set(CMAKE_CXX_STANDARD_REQUIRED ON)
set(CMAKE_EXPORT_COMPILE_COMMANDS ON)
set(CMAKE_RUNTIME_OUTPUT_DIRECTORY ${CMAKE_BINARY_DIR}/bin)

if(NOT CMAKE_BUILD_TYPE AND NOT CMAKE_CONFIGURATION_TYPES)
    set(CMAKE_BUILD_TYPE "RelWithDebInfo" CACHE STRING "" FORCE)
endif()

add_subdirectory(libs/math)
add_subdirectory(app)

option(BUILD_TESTS "Build the unit tests" ON)
if(BUILD_TESTS)
    enable_testing()
    add_subdirectory(tests)
endif()
\end{lstlisting}

\begin{lstlisting}[title={\color{ortaGri}\ttfamily libs/math/CMakeLists.txt}]
add_library(math STATIC src/calculator.cpp)
add_library(Calc::math ALIAS math)

target_include_directories(math PUBLIC
    $<BUILD_INTERFACE:${CMAKE_CURRENT_SOURCE_DIR}/include>
    $<INSTALL_INTERFACE:include>
)
target_compile_features(math PUBLIC cxx_std_17)
\end{lstlisting}

\begin{lstlisting}[title={\color{ortaGri}\ttfamily app/CMakeLists.txt}]
add_executable(calc_app main.cpp)
target_link_libraries(calc_app PRIVATE Calc::math)
\end{lstlisting}

\begin{lstlisting}[style=cppstyle,title={\color{ortaGri}\ttfamily app/main.cpp}]
#include <iostream>
#include "math/calculator.h"

int main() {
    std::cout << "2 + 3   = " << calc::add(2, 3)      << "\n";
    std::cout << "4 * 5   = " << calc::multiply(4, 5) << "\n";
    std::cout << "5!      = " << calc::factorial(5)   << "\n";
    return 0;
}
\end{lstlisting}

\begin{lstlisting}[title={\color{ortaGri}\ttfamily tests/CMakeLists.txt}]
add_executable(calc_test calc_test.cpp)
target_link_libraries(calc_test PRIVATE Calc::math)
add_test(NAME CalcBasic COMMAND calc_test)
\end{lstlisting}

\begin{lstlisting}[style=cppstyle,title={\color{ortaGri}\ttfamily tests/calc\_test.cpp}]
#include <cassert>
#include "math/calculator.h"

int main() {
    assert(calc::add(2, 3) == 5.0);
    assert(calc::multiply(4, 5) == 20.0);
    assert(calc::factorial(5) == 120);
    return 0;   // if every assert passes, 0 = success
}
\end{lstlisting}

\gsub{Derleme ve Çalıştırma}

\begin{lstlisting}[style=shstyle]
cmake -S . -B build -G Ninja
cmake --build build
./build/bin/calc_app
#   2 + 3   = 5
#   4 * 5   = 20
#   5!      = 120
cd build && ctest --output-on-failure
#   100% tests passed, 0 tests failed out of 1
\end{lstlisting}

% =====================================================================
\gsec{Tam Örnek: C Projesi (Kütüphane + Uygulama)}
% =====================================================================

Aynı yapıyı bu kez saf C ile kuralım: bir dinamik dizi (dynamic array)
kütüphanesi ve onu kullanan bir uygulama.

\gsub{Dosya Yapısı}

\begin{lstlisting}[style=shstyle]
darray/
|-- CMakeLists.txt
|-- lib/
|   |-- CMakeLists.txt
|   |-- include/darray.h
|   `-- src/darray.c
`-- app/
    |-- CMakeLists.txt
    `-- main.c
\end{lstlisting}

\gsub{Kaynaklar}

\begin{lstlisting}[style=cstyle,title={\color{ortaGri}\ttfamily lib/include/darray.h}]
#ifndef DARRAY_H
#define DARRAY_H
#include <stddef.h>

typedef struct {
    int   *data;
    size_t size;
    size_t capacity;
} DArray;

void darray_init(DArray *a);
void darray_push(DArray *a, int value);
void darray_free(DArray *a);

#endif /* DARRAY_H */
\end{lstlisting}

\begin{lstlisting}[style=cstyle,title={\color{ortaGri}\ttfamily lib/src/darray.c}]
#include "darray.h"
#include <stdlib.h>

void darray_init(DArray *a) {
    a->data = NULL; a->size = 0; a->capacity = 0;
}

void darray_push(DArray *a, int value) {
    if (a->size == a->capacity) {
        a->capacity = a->capacity ? a->capacity * 2 : 4;
        a->data = realloc(a->data, a->capacity * sizeof(int));
    }
    a->data[a->size++] = value;
}

void darray_free(DArray *a) {
    free(a->data);
    a->data = NULL; a->size = a->capacity = 0;
}
\end{lstlisting}

\begin{lstlisting}[style=cstyle,title={\color{ortaGri}\ttfamily app/main.c}]
#include <stdio.h>
#include "darray.h"

int main(void) {
    DArray a;
    darray_init(&a);
    for (int i = 1; i <= 5; ++i) darray_push(&a, i * i);

    printf("Size: %zu\n", a.size);
    for (size_t i = 0; i < a.size; ++i) printf("%d ", a.data[i]);
    printf("\n");

    darray_free(&a);
    return 0;
}
\end{lstlisting}

\gsub{CMakeLists.txt Dosyaları}

\begin{lstlisting}[title={\color{ortaGri}\ttfamily CMakeLists.txt (kök)}]
cmake_minimum_required(VERSION 3.20)
project(DArray VERSION 1.0 LANGUAGES C)

set(CMAKE_C_STANDARD 11)
set(CMAKE_C_STANDARD_REQUIRED ON)
set(CMAKE_RUNTIME_OUTPUT_DIRECTORY ${CMAKE_BINARY_DIR}/bin)

add_subdirectory(lib)
add_subdirectory(app)
\end{lstlisting}

\begin{lstlisting}[title={\color{ortaGri}\ttfamily lib/CMakeLists.txt}]
add_library(darray STATIC src/darray.c)
add_library(DArray::darray ALIAS darray)

target_include_directories(darray PUBLIC
    ${CMAKE_CURRENT_SOURCE_DIR}/include
)
target_compile_features(darray PUBLIC c_std_11)

# Strict warnings only on GCC/Clang (MSVC would reject these flags)
target_compile_options(darray PRIVATE
    $<$<C_COMPILER_ID:GNU,Clang>:-Wall;-Wextra>
)
\end{lstlisting}

\begin{lstlisting}[title={\color{ortaGri}\ttfamily app/CMakeLists.txt}]
add_executable(darray_app main.c)
target_link_libraries(darray_app PRIVATE DArray::darray)
\end{lstlisting}

\begin{lstlisting}[style=shstyle]
cmake -S . -B build -G Ninja
cmake --build build
./build/bin/darray_app
#   Size: 5
#   1 4 9 16 25
\end{lstlisting}

% =====================================================================
\gsec{En İyi Uygulamalar}
% =====================================================================

\begin{center}
\begin{tabularx}{\textwidth}{@{}>{\bfseries\color{yesil}}c >{\raggedright\arraybackslash}X@{}}
\toprule
\rowcolor{acikYesil}\color{yesil}\textbf{Yap} & \color{yesil}\textbf{Neden}\\
\midrule
\checkmark & \code{target\_*} komutlarını kullan; global \code{include\_directories}/\code{add\_definitions}'dan kaçın.\\
\rowcolor{acikGri}\color{yesil}\checkmark & PUBLIC/PRIVATE/INTERFACE'i bilinçli seç --- kullanım gereksinimleri kendiliğinden yayılsın.\\
\checkmark & Out-of-source build yap (\code{-B build}); \code{build/} klasörünü git'e ekleme.\\
\rowcolor{acikGri}\color{yesil}\checkmark & Kütüphanelere \code{Ad::Alt} ALIAS ver; \code{find\_package} hedefleriyle tutarlı olsun.\\
\checkmark & \code{FetchContent}'te daima sabit \code{GIT\_TAG} (sürüm/commit) kullan.\\
\rowcolor{acikGri}\color{yesil}\checkmark & \code{cmake\_minimum\_required} sürümünü gerçekçi ve güncel tut (3.20+).\\
\checkmark & \code{CMAKE\_EXPORT\_COMPILE\_COMMANDS ON} ile IDE/clangd desteğini aç.\\
\rowcolor{acikGri}\color{yesil}\checkmark & Derleyiciye bağlı bayrakları genex (\code{\$<...>}) ile koşulla.\\
\bottomrule
\end{tabularx}
\end{center}

\begin{center}
\begin{tabularx}{\textwidth}{@{}>{\bfseries\color{kirmizi}}c >{\raggedright\arraybackslash}X@{}}
\toprule
\rowcolor{kirmizi!8}\color{kirmizi}\textbf{Yapma} & \color{kirmizi}\textbf{Neden}\\
\midrule
\texttimes & \code{file(GLOB ...)} ile kaynak toplama --- yeni dosya eklenince CMake yeniden çalışmaz, derleme kırılır.\\
\rowcolor{acikGri}\color{kirmizi}\texttimes & Kaynak dizinine derleme yapma (in-source build) --- kaynağı kirletir.\\
\texttimes & \code{link\_libraries}/\code{include\_directories}'i global kullanma.\\
\rowcolor{acikGri}\color{kirmizi}\texttimes & Sabit yollar (\code{/home/tolga/...}) yazma --- taşınabilir değildir.\\
\texttimes & \code{CMAKE\_CXX\_FLAGS}'ı elle üzerine yazma --- mevcut bayrakları ezersin.\\
\bottomrule
\end{tabularx}
\end{center}

\uyari{\code{file(GLOB SOURCES "src/*.cpp")} cazip görünür ama tehlikelidir:
CMake dosya listesini yapılandırma anında dondurur; yeni bir \code{.cpp}
eklediğinizde CMake bunu fark etmez ve dosya derlemeye girmez. Kaynakları
her zaman \code{add\_executable}/\code{add\_library} içinde \emph{elle}
listeleyin. (3.12+ ile \code{CONFIGURE\_DEPENDS} kısmen yardımcı olsa da yine
de önerilmez.)}

% =====================================================================
\gsec{Sık Karşılaşılan Hatalar ve Çözümleri}
% =====================================================================

\begin{center}
\begin{tabularx}{\textwidth}{@{}>{\raggedright\arraybackslash}X >{\raggedright\arraybackslash}X@{}}
\toprule
\rowcolor{acikMavi}\color{anaMavi}\textbf{Hata / Belirti} & \color{anaMavi}\textbf{Neden \& Çözüm}\\
\midrule
\code{undefined reference to `foo'} & Bağlayıcı sembolü bulamıyor. Fonksiyonu tanımlayan kütüphaneyi \code{target\_link\_libraries} ile ekleyin.\\
\rowcolor{acikGri} \code{fatal error: xyz.h: No such file} & Include yolu eksik. \code{target\_include\_directories} ile başlık dizinini ekleyin (doğru PUBLIC/PRIVATE ile).\\
\code{Could not find package XYZ} & \code{find\_package} kütüphaneyi bulamadı. \code{-DCMAKE\_PREFIX\_PATH=...} verin veya kütüphaneyi kurun.\\
\rowcolor{acikGri} \code{CMake Error: ... does not appear to contain CMakeLists.txt} & \code{-S} yanlış dizini gösteriyor. Kök \code{CMakeLists.txt}'nin bulunduğu dizini verin.\\
Değişiklik derlemeye yansımıyor & Eski cache. \code{build/} klasörünü silin ya da \code{cmake -{}-fresh} ile yeniden yapılandırın.\\
\rowcolor{acikGri} \code{target\_link\_libraries ... is not a target} & Hedef henüz tanımlanmamış. \code{add\_subdirectory} sırasını düzeltin (kütüphane önce).\\
Release'te optimizasyon yok & \code{CMAKE\_BUILD\_TYPE} boş. \code{-DCMAKE\_BUILD\_TYPE=Release} verin.\\
\bottomrule
\end{tabularx}
\end{center}

% =====================================================================
\gsec{Hızlı Başvuru (Cheat Sheet)}
% =====================================================================

\gsub{En Sık Komutlar}

\begin{center}
\begin{tabularx}{\textwidth}{@{}>{\ttfamily\footnotesize\color{koyuGri}}l >{\raggedright\arraybackslash}X@{}}
\toprule
\rowcolor{acikMavi}\color{anaMavi}\textbf{Komut} & \color{anaMavi}\textbf{Ne yapar}\\
\midrule
cmake\_minimum\_required(VERSION x) & Gereken en düşük CMake sürümü (ilk satır).\\
\rowcolor{acikGri} project(Ad VERSION x LANGUAGES CXX C) & Projeyi ve dilleri tanımlar.\\
add\_executable(ad kaynaklar) & Çalıştırılabilir hedef.\\
\rowcolor{acikGri} add\_library(ad STATIC/SHARED src) & Kütüphane hedefi.\\
target\_link\_libraries(t PUBLIC x) & Hedefe kütüphane bağlar (özellikler yayılır).\\
\rowcolor{acikGri} target\_include\_directories(t PUBLIC d) & Başlık arama yolu.\\
target\_compile\_features(t PUBLIC cxx\_std\_17) & Dil standardı gereksinimi.\\
\rowcolor{acikGri} target\_compile\_definitions(t PRIVATE M) & Önişlemci makrosu.\\
find\_package(X REQUIRED) & Harici kütüphane bulur.\\
\rowcolor{acikGri} add\_subdirectory(dir) & Alt dizindeki CMakeLists'i dahil eder.\\
install(TARGETS ...) & Kurulum kuralları.\\
\rowcolor{acikGri} enable\_testing() / add\_test(...) & CTest testleri.\\
\bottomrule
\end{tabularx}
\end{center}

\gsub{En Sık Terminal Komutları}

\begin{lstlisting}[style=shstyle]
# Configure + generate (Ninja, Release)
cmake -S . -B build -G Ninja -D CMAKE_BUILD_TYPE=Release

# Build (with 8 parallel jobs)
cmake --build build -j 8

# Build a single target only
cmake --build build --target calc_app

# Run the tests
ctest --test-dir build --output-on-failure

# Install
cmake --install build --prefix /opt/calc

# Reconfigure from scratch (wipe the cache)
cmake --fresh -S . -B build
# or:  rm -rf build

# Using a preset
cmake --preset debug && cmake --build --preset debug
\end{lstlisting}

\gsub{Önemli Yerleşik Değişkenler}

\begin{center}
\begin{tabularx}{\textwidth}{@{}>{\ttfamily\footnotesize\color{koyuGri}}l >{\raggedright\arraybackslash}X@{}}
\toprule
\rowcolor{acikMavi}\color{anaMavi}\textbf{Değişken} & \color{anaMavi}\textbf{Anlamı}\\
\midrule
CMAKE\_SOURCE\_DIR & En üst kaynak dizini.\\
\rowcolor{acikGri} CMAKE\_BINARY\_DIR & En üst derleme dizini.\\
CMAKE\_CURRENT\_SOURCE\_DIR & İçinde bulunulan CMakeLists'in dizini.\\
\rowcolor{acikGri} PROJECT\_NAME / PROJECT\_VERSION & En yakın project() değerleri.\\
CMAKE\_BUILD\_TYPE & Debug/Release/... (tek-config generator).\\
\rowcolor{acikGri} CMAKE\_INSTALL\_PREFIX & Kurulum kök dizini.\\
BUILD\_SHARED\_LIBS & Tip belirtilmeyen kütüphaneler SHARED mı?\\
\rowcolor{acikGri} CMAKE\_CXX\_STANDARD & Varsayılan C++ standardı.\\
\bottomrule
\end{tabularx}
\end{center}

\ipucu{Bu rehber seni derleme sistemlerinin \emph{neden} var olduğundan
başlayıp, CMake'in hedef-tabanlı modern stiliyle çok modüllü, testli,
paketlenebilir projeler kurmaya taşıdı: derleyici $\to$ make $\to$ CMake
$\to$ hedefler $\to$ bağımlılıklar $\to$ test $\to$ paketleme. Buradan sonrası
pratik: küçük bir kütüphane yaz, onu \code{find\_package}'le tükettir,
CI'a bağla. Her \code{CMakeLists.txt} satırını ``bu ayar hangi hedefe ait?''
sorusuyla oku --- modern CMake'in tamamı bu sorunun cevabıdır. İyi derlemeler!}

\vfill
\begin{center}
{\color{ortaGri}\rule{0.5\linewidth}{0.6pt}}\\[6pt]
{\color{ortaGri}\small CMake ile Modern C/C++ Derleme Rehberi --- Son}
\end{center}

\end{document}
